\documentclass[8pt]{beamer}
\usetheme{Madrid}
\usecolortheme{default}
\usepackage{amsmath, amssymb, graphicx, array, multirow, booktabs, xcolor, tikz}
\usetikzlibrary{shapes,arrows,positioning,fit,backgrounds}

% Compact font settings
\setbeamerfont{title}{size=\Large}
\setbeamerfont{subtitle}{size=\small}
\setbeamerfont{author}{size=\scriptsize}
\setbeamerfont{date}{size=\scriptsize}
\setbeamerfont{frametitle}{size=\large}
\setbeamerfont{framesubtitle}{size=\normalsize}
\setbeamerfont{enumerate item}{size=\footnotesize}
\setbeamerfont{itemize item}{size=\footnotesize}
\setbeamerfont{block title}{size=\footnotesize}

% Compact spacing
\setlength{\itemsep}{0.08ex}
\setlength{\parskip}{0.08ex}
\setlength{\leftmargini}{0.3cm}
\setbeamersize{text margin left=3mm, text margin right=3mm}
\addtobeamertemplate{frame begin}{\vspace{-0.4cm}}{}

% Colors
\definecolor{myblue}{RGB}{0,0,255}
\definecolor{myred}{RGB}{255,0,0}
\definecolor{mygreen}{RGB}{0,128,0}
\definecolor{myorange}{RGB}{255,165,0}
\definecolor{mypurple}{RGB}{128,0,128}
\definecolor{mygray}{RGB}{128,128,128}
\definecolor{mygold}{RGB}{255,215,0}

\title{Generative Artificial Intelligence}
\subtitle{Complete Tutorial: GenAI, LLMs, Transformers, Word2Vec, RNNs, LSTM}
\author{GARP RAI Exam Preparation}
\date{}

\begin{document}
	
	\begin{frame}
		\titlepage
	\end{frame}
	
	% ============================================================
	% SECTION 1: INTRODUCTION TO GENAI - SLIDES 1-10
	% ============================================================
	
	\begin{frame}{Slide 1: What is Generative AI?}
		\fontsize{6.5}{7.5}\selectfont
		\textbf{Definition:}
		\begin{itemize}
			\item Generative Artificial Intelligence (GenAI) is a branch of AI that employs models to generate text, sound, or videos
			\item Creates new content rather than just analyzing existing data
			\item Represents one of the most significant technological developments of the 2020s
		\end{itemize}
		
		\vspace{0.1cm}
		\textbf{Key Milestone:}
		\begin{itemize}
			\item OpenAI released ChatGPT in November 2022
			\item Reached 1 million users in just 5 days (fastest AI tool ever)
			\item Facebook took almost a year to reach this milestone
		\end{itemize}
		
		\vspace{0.1cm}
		\textbf{Impact:}
		\begin{itemize}
			\item Facilitated rapid adoption of Generative AI
			\item Transformed how businesses and individuals interact with AI
			\item Created new opportunities across all sectors
		\end{itemize}
		
		\textcolor{myblue}{\textbf{Key Insight:}} GenAI creates new content - not just analyzes!
	\end{frame}
	
	\begin{frame}{Slide 2: GenAI vs Traditional AI}
		\fontsize{6.5}{7.5}\selectfont
		\textbf{Traditional AI/Machine Learning:}
		\begin{itemize}
			\item Analyzes existing data
			\item Makes predictions or classifications
			\item Examples: Regression, classification, clustering
			\item Outputs are labels or numerical values
			\item Deterministic (same input = same output)
		\end{itemize}
		
		\vspace{0.1cm}
		\textbf{Generative AI:}
		\begin{itemize}
			\item Creates new content
			\item Generates text, images, audio, video
			\item Examples: ChatGPT, DALL-E, Gemini
			\item Outputs are novel, creative content
			\item Probabilistic (same input may yield different outputs)
		\end{itemize}
		
		\vspace{0.1cm}
		\textbf{Key Difference:}
		\begin{itemize}
			\item Traditional AI: "What is this?"
			\item Generative AI: "Create something new based on this"
		\end{itemize}
		
		\textcolor{myblue}{\textbf{Remember:}} GenAI creates; Traditional AI analyzes!
	\end{frame}
	
	\begin{frame}{Slide 3: ChatGPT Overview}
		\fontsize{6.5}{7.5}\selectfont
		\textbf{What is ChatGPT?}
		\begin{itemize}
			\item Chatbot based on GPT (Generative Pre-trained Transformer)
			\item Developed by OpenAI
			\item Uses transformer-based LLM technology
			\item Mimics human conversation via text
		\end{itemize}
		
		\vspace{0.1cm}
		\textbf{Key Features:}
		\begin{enumerate}
			\item Conversational interface
			\item Can answer follow-up questions
			\item Admits mistakes
			\item Challenges incorrect premises
			\item Rejects inappropriate requests
		\end{enumerate}
		
		\vspace{0.1cm}
		\textbf{Significance:}
		\begin{itemize}
			\item Made LLMs accessible to general public
			\item Demonstrated practical applications of GenAI
			\item Sparked widespread interest and adoption
		\end{itemize}
		
		\textcolor{myblue}{\textbf{Key Insight:}} ChatGPT popularized LLMs for everyday use!
	\end{frame}
	
	\begin{frame}{Slide 4: GenAI Taxonomy - Content Modalities}
		\fontsize{6.5}{7.5}\selectfont
		\textbf{Five Major Content Modalities:}
		\begin{enumerate}
			\item \textbf{Text}
			\begin{itemize}
				\item Underlying technology: Transformers
				\item Examples: ChatGPT, Claude, Gemini
			\end{itemize}
			\item \textbf{Image}
			\begin{itemize}
				\item Underlying technology: DALL-E, Stable Diffusion
				\item Creates images from text prompts
			\end{itemize}
			\item \textbf{Audio}
			\begin{itemize}
				\item Includes speech, sound effects, music
				\item Examples: Text-to-speech, music generation
			\end{itemize}
			\item \textbf{Video}
			\begin{itemize}
				\item Creates videos from prompts
				\item Examples: Sora, RunwayML
			\end{itemize}
			\item \textbf{Multimodal}
			\begin{itemize}
				\item Takes multiple input types
				\item Generates multiple output types
				\item Examples: Gemini, GPT-4
			\end{itemize}
		\end{enumerate}
		
		\textcolor{myblue}{\textbf{Key Concept:}} GenAI works across multiple content types!
	\end{frame}
	
	\begin{frame}{Slide 5: GenAI Taxonomy - Modeling Algorithms}
		\fontsize{6.5}{7.5}\selectfont
		\textbf{Key Modeling Algorithms:}
		\begin{enumerate}
			\item \textbf{Large Language Models (LLMs)}
			\begin{itemize}
				\item Billions of parameters
				\item Trained on vast text corpora
				\item Transformer-based architecture
			\end{itemize}
			\item \textbf{Transformers}
			\begin{itemize}
				\item Self-attention mechanism
				\item Examples: BERT, GPT
			\end{itemize}
			\item \textbf{Stable Diffusion}
			\begin{itemize}
				\item Image generation
				\item Probabilistic reconstruction
			\end{itemize}
			\item \textbf{Variational Autoencoders (VAE)}
			\begin{itemize}
				\item Deep neural networks
				\item Image modification and generation
			\end{itemize}
			\item \textbf{Generative Adversarial Networks (GANs)}
			\begin{itemize}
				\item Generator + Discriminator
				\item Creates realistic synthetic data
			\end{itemize}
			\item \textbf{Recurrent Neural Networks (RNNs)}
			\begin{itemize}
				\item Sequential data processing
				\item Time series and text
			\end{itemize}
		\end{enumerate}
		
		\textcolor{myblue}{\textbf{Key Insight:}} Different algorithms for different content types!
	\end{frame}
	
	\begin{frame}{Slide 6: DALL-E and Image Generation}
		\fontsize{6.5}{7.5}\selectfont
		\textbf{DALL-E:}
		\begin{itemize}
			\item Popular image generator developed by OpenAI
			\item Takes text prompts and creates images
			\item Can generate novel, creative images
		\end{itemize}
		
		\vspace{0.1cm}
		\textbf{Example Prompt:}
		\begin{itemize}
			\item "Create an image that personifies the 'ups and downs' of financial markets as if it were investors sitting on a roller coaster"
		\end{itemize}
		
		\vspace{0.1cm}
		\textbf{Capabilities:}
		\begin{enumerate}
			\item Generate images from text descriptions
			\item Modify existing images
			\item Combine concepts in creative ways
			\item Understand complex prompts
		\end{enumerate}
		
		\vspace{0.1cm}
		\textbf{Applications:}
		\begin{itemize}
			\item Marketing materials
			\item Rapid prototyping
			\item Creative projects
			\item Presentations and illustrations
		\end{itemize}
		
		\textcolor{myblue}{\textbf{Key Insight:}} DALL-E creates images from text descriptions!
	\end{frame}
	
	\begin{frame}{Slide 7: Stable Diffusion Model}
		\fontsize{6.5}{7.5}\selectfont
		\textbf{What is Stable Diffusion?}
		\begin{itemize}
			\item Developed by Stability AI (released 2022)
			\item Powerful image generating technology
			\item Trained on database of images
			\item Uses probabilistic functions
		\end{itemize}
		
		\vspace{0.1cm}
		\textbf{How It Works:}
		\begin{enumerate}
			\item Start with random noise
			\item Gradually refine to match text prompt
			\item Uses diffusion process
			\item Reconstructs images step by step
		\end{enumerate}
		
		\vspace{0.1cm}
		\textbf{Advantages:}
		\begin{itemize}
			\item Open source
			\item Runs on consumer GPUs
			\item High quality outputs
			\item Customizable
		\end{itemize}
		
		\textcolor{myblue}{\textbf{Key Concept:}} Stable Diffusion reconstructs images from noise based on text!
	\end{frame}
	
	\begin{frame}{Slide 8: Variational Autoencoders (VAE)}
		\fontsize{6.5}{7.5}\selectfont
		\textbf{What are VAEs?}
		\begin{itemize}
			\item Type of deep neural network
			\item Multiple hidden layers
			\item Used for image modality
			\item Particularly useful for modifying existing images
		\end{itemize}
		
		\vspace{0.1cm}
		\textbf{Architecture:}
		\begin{enumerate}
			\item \textbf{Encoding Layer}
			\begin{itemize}
				\item Compresses images
				\item Reduces dimensionality
			\end{itemize}
			\item \textbf{Decoding Layer}
			\begin{itemize}
				\item Reconstructs images
				\item Uses probabilistic inference
				\item Fills in missing information
			\end{itemize}
		\end{enumerate}
		
		\vspace{0.1cm}
		\textbf{Applications:}
		\begin{itemize}
			\item Image modification
			\item Anomaly detection
			\item Data generation
			\item Feature extraction
		\end{itemize}
		
		\textcolor{myblue}{\textbf{Key Insight:}} VAEs are excellent for modifying existing images!
	\end{frame}
	
	\begin{frame}{Slide 9: Generative Adversarial Networks (GANs)}
		\fontsize{6.5}{7.5}\selectfont
		\textbf{What are GANs?}
		\begin{itemize}
			\item Used to create images, audio, and video
			\item Employs two competing neural networks
			\item Generator creates simulated data
			\item Discriminator distinguishes real from fake
		\end{itemize}
		
		\vspace{0.1cm}
		\textbf{Architecture:}
		\begin{enumerate}
			\item \textbf{Generator}
			\begin{itemize}
				\item Creates simulated data
				\item Learns to fool discriminator
			\end{itemize}
			\item \textbf{Discriminator}
			\begin{itemize}
				\item Identifies real vs fake data
				\item Assigns probability scores
				\item Sends feedback to generator
			\end{itemize}
		\end{enumerate}
		
		\vspace{0.1cm}
		\textbf{Process:}
		\begin{itemize}
			\item Iterative competition
			\item Generator improves over time
			\item Eventually generates indistinguishable data
		\end{itemize}
		
		\textcolor{myblue}{\textbf{Key Insight:}} GANs = two networks competing to create realistic content!
	\end{frame}
	
	\begin{frame}{Slide 10: GenAI vs LLM Relationship}
		\fontsize{6.5}{7.5}\selectfont
		\textbf{Relationship Overview:}
		\begin{itemize}
			\item LLMs are a subset of GenAI
			\item GenAI is broader - includes image, audio, video
			\item LLMs specifically focus on text generation
		\end{itemize}
		
		\vspace{0.1cm}
		\textbf{GenAI Includes:}
		\begin{enumerate}
			\item \textbf{Text Models (LLMs)}
			\begin{itemize}
				\item ChatGPT, Claude, Gemini
			\end{itemize}
			\item \textbf{Image Models}
			\begin{itemize}
				\item DALL-E, Stable Diffusion
			\end{itemize}
			\item \textbf{Audio Models}
			\begin{itemize}
				\item Text-to-speech, music generation
			\end{itemize}
			\item \textbf{Video Models}
			\begin{itemize}
				\item Sora, RunwayML
			\end{itemize}
			\item \textbf{Multimodal Models}
			\begin{itemize}
				\item Gemini, GPT-4
			\end{itemize}
		\end{enumerate}
		
		\textcolor{myblue}{\textbf{Key Concept:}} All LLMs are GenAI, but not all GenAI are LLMs!
	\end{frame}
	
	% ============================================================
	% SECTION 2: WORD EMBEDDINGS & WORD2VEC - SLIDES 11-25
	% ============================================================
	
	\begin{frame}{Slide 11: BoW Limitations Review}
		\fontsize{6.5}{7.5}\selectfont
		\textbf{Three Major Limitations:}
		\begin{enumerate}
			\item \textbf{No Semantic Understanding}
			\begin{itemize}
				\item Cannot identify similar meanings
				\item "Amazing" and "fantastic" treated differently
			\end{itemize}
			\item \textbf{Sparse Vectors}
			\begin{itemize}
				\item Vast vectors with mostly zeros
				\item Inefficient storage and analysis
				\item Computational burden
			\end{itemize}
			\item \textbf{No Context}
			\begin{itemize}
				\item Words treated independently
				\item "Call option" vs "phone call" confused
				\item Negation words not handled
			\end{itemize}
		\end{enumerate}
		
		\vspace{0.1cm}
		\textbf{Why This Matters:}
		\begin{itemize}
			\item Limits analytical power
			\item Misses important linguistic relationships
			\item Creates computational inefficiency
		\end{itemize}
		
		\textcolor{myblue}{\textbf{Key Insight:}} BoW has fundamental limitations Word2Vec addresses!
	\end{frame}
	
	\begin{frame}{Slide 12: Word2Vec Introduction}
		\fontsize{6.5}{7.5}\selectfont
		\textbf{What is Word2Vec?}
		\begin{itemize}
			\item NLP model created by Google team (2013)
			\item Led by Tomas Mikolov
			\item Alternative to conventional BoW
			\item Represents text as numerical vectors
		\end{itemize}
		
		\vspace{0.1cm}
		\textbf{Key Concept:}
		\begin{itemize}
			\item Each word has its own vector
			\item Constructed by examining surrounding words
			\item Measures "co-occurrence" of words
		\end{itemize}
		
		\vspace{0.1cm}
		\textbf{Advantages Over BoW:}
		\begin{enumerate}
			\item Compressed representation
			\item Dense vector representations
			\item Reduces dimensionality
			\item Captures semantic relationships
		\end{enumerate}
		
		\textcolor{myblue}{\textbf{Key Insight:}} Word2Vec creates dense, meaningful word representations!
	\end{frame}
	
	\begin{frame}{Slide 13: Word Embeddings Explained}
		\fontsize{6.5}{7.5}\selectfont
		\textbf{What are Word Embeddings?}
		\begin{itemize}
			\item Dense vector representations of words
			\item Generated through neural network training
			\item Encode semantic information
			\item Position words in continuous vector space
		\end{itemize}
		
		\vspace{0.1cm}
		\textbf{Characteristics:}
		\begin{enumerate}
			\item \textbf{Dense:} Fewer dimensions, mostly non-zero values
			\item \textbf{Semantic:} Similar words have similar vectors
			\item \textbf{Learnable:} Trained from data
			\item \textbf{Contextual:} Capture word relationships
		\end{enumerate}
		
		\vspace{0.1cm}
		\textbf{Typical Dimensions:}
		\begin{itemize}
			\item Embedding layer often has 300 neurons
			\item Much smaller than vocabulary size
			\item Efficient representation
		\end{itemize}
		
		\textcolor{myblue}{\textbf{Key Concept:}} Embeddings = dense vectors capturing meaning!
	\end{frame}
	
	\begin{frame}{Slide 14: Continuous Bag of Words (CBoW)}
		\fontsize{6.5}{7.5}\selectfont
		\textbf{CBoW Architecture:}
		\begin{itemize}
			\item "Fill in the blank" technique
			\item Predicts missing word from context
			\item Uses context words before and after
		\end{itemize}
		
		\vspace{0.1cm}
		\textbf{How It Works:}
		\begin{enumerate}
			\item Set context window (size c)
			\item Identify words before and after target (2c words)
			\item Input: context words
			\item Output: target word
			\item Goal: $P(w_i | w_{i-c}, ..., w_{i+c})$
		\end{enumerate}
		
		\vspace{0.1cm}
		\textbf{Example:}
		\begin{itemize}
			\item Sentence: "bonds are less risky compared to stocks"
			\item Target word: "risky"
			\item Input: words preceding and following "risky"
		\end{itemize}
		
		\textcolor{myblue}{\textbf{Key Insight:}} CBoW predicts target from surrounding context!
	\end{frame}
	
	\begin{frame}{Slide 15: Skip-gram Architecture}
		\fontsize{6.5}{7.5}\selectfont
		\textbf{Skip-gram Architecture:}
		\begin{itemize}
			\item "Word association" technique
			\item Uses one word to predict surrounding words
			\item Maximizes joint probability of context words
		\end{itemize}
		
		\vspace{0.1cm}
		\textbf{How It Works:}
		\begin{enumerate}
			\item Input: target word
			\item Output: context words (2c words)
			\item Goal: $P(w_{i-c}, ..., w_{i+c} | w_i)$
			\item Tries to predict surrounding words
		\end{enumerate}
		
		\vspace{0.1cm}
		\textbf{Example:}
		\begin{itemize}
			\item Sentence: "bonds are less risky compared to stocks"
			\item Input: "risky"
			\item Output: words preceding and following "risky"
		\end{itemize}
		
		\vspace{0.1cm}
		\textbf{Comparison with CBoW:}
		\begin{itemize}
			\item CBoW: context $\to$ target
			\item Skip-gram: target $\to$ context
		\end{itemize}
		
		\textcolor{myblue}{\textbf{Key Insight:}} Skip-gram predicts context from target word!
	\end{frame}
	
	\begin{frame}{Slide 16: Word2Vec Neural Network Architecture}
		\fontsize{6.5}{7.5}\selectfont
		\textbf{Network Structure:}
		\begin{itemize}
			\item Similar to autoencoder
			\item Single hidden layer (embedding layer)
			\item Input: one-hot encoded vectors
			\item Output: probability predictions
		\end{itemize}
		
		\vspace{0.1cm}
		\textbf{Components:}
		\begin{enumerate}
			\item \textbf{Input Layer}
			\begin{itemize}
				\item One-hot encoded vectors
				\item Size: vocabulary size ($|W|$)
			\end{itemize}
			\item \textbf{Hidden Layer}
			\begin{itemize}
				\item Embedding layer
				\item Typically 300 neurons
				\item Learns word representations
			\end{itemize}
			\item \textbf{Output Layer}
			\begin{itemize}
				\item Probability distribution over words
				\item Size: vocabulary size ($|W|$)
			\end{itemize}
		\end{enumerate}
		
		\textcolor{myblue}{\textbf{Key Insight:}} Hidden layer creates the word embeddings!
	\end{frame}
	
	\begin{frame}{Slide 17: Dimensionality Reduction with Word2Vec}
		\fontsize{6.5}{7.5}\selectfont
		\textbf{Comparison of Dimensions:}
		
		\begin{tabular}{|l|c|}
			\hline
			\textbf{Method} & \textbf{Dimensions} \\
			\hline
			Traditional BoW & $|W| \times |W|$ \\
			Word2Vec & $|W| \times 300$ \\
			\hline
		\end{tabular}
		
		\vspace{0.1cm}
		\textbf{Example Calculation:}
		\begin{itemize}
			\item Vocabulary size: 100,000 words
			\item Traditional BoW: 100,000 $\times$ 100,000 = 10 billion weights
			\item Word2Vec: 100,000 $\times$ 300 = 30 million weights
			\item \textbf{Reduction:} 99.7\% fewer weights!
		\end{itemize}
		
		\vspace{0.1cm}
		\textbf{Benefits of Reduction:}
		\begin{enumerate}
			\item Less memory required
			\item Faster computation
			\item Better generalization
			\item More efficient training
		\end{enumerate}
		
		\textcolor{myblue}{\textbf{Key Insight:}} Word2Vec dramatically reduces dimensionality!
	\end{frame}
	
	\begin{frame}{Slide 18: Training Word2Vec}
		\fontsize{6.5}{7.5}\selectfont
		\textbf{Training Process:}
		\begin{enumerate}
			\item \textbf{Feed neural network with text}
			\begin{itemize}
				\item Large corpus of documents
				\item Process word by word
			\end{itemize}
			\item \textbf{Learn embedding weights}
			\begin{itemize}
				\item Through backpropagation
				\item Using gradient descent
			\end{itemize}
			\item \textbf{Generate dense vectors}
			\begin{itemize}
				\item Each word gets vector representation
				\item Semantic information encoded
			\end{itemize}
		\end{enumerate}
		
		\vspace{0.1cm}
		\textbf{What the Model Learns:}
		\begin{itemize}
			\item Co-occurrence patterns
			\item Word relationships
			\item Semantic meanings
			\item Contextual information
		\end{itemize}
		
		\textcolor{myblue}{\textbf{Key Concept:}} Training learns embeddings from text data!
	\end{frame}
	
	\begin{frame}{Slide 19: Cosine Similarity}
		\fontsize{6.5}{7.5}\selectfont
		\textbf{Definition:}
		\begin{itemize}
			\item Measures "distance" between words in embedding space
			\item Similar words close together
			\item Dissimilar words far apart
			\item Range: [-1, 1]
		\end{itemize}
		
		\vspace{0.1cm}
		\textbf{Formula:}
		\begin{itemize}
			\item $\cos(\theta) = \frac{P \cdot Q}{||P|| \times ||Q||}$
		\end{itemize}
		
		\vspace{0.1cm}
		\textbf{Interpretation:}
		\begin{tabular}{|c|c|}
			\hline
			\textbf{Cosine Similarity} & \textbf{Meaning} \\
			\hline
			1.0 & Identical (or very similar) \\
			0.7 - 0.9 & Very similar \\
			0.4 - 0.6 & Somewhat similar \\
			0.0 - 0.3 & Dissimilar \\
			\hline
		\end{tabular}
		
		\vspace{0.1cm}
		\textbf{Example:}
		\begin{itemize}
			\item China and Japan: 0.84
			\item Gold and Silver: 0.95
			\item Banana and Rose: 0.30
		\end{itemize}
		
		\textcolor{myblue}{\textbf{Key Insight:}} Cosine similarity measures word closeness!
	\end{frame}
	
	\begin{frame}{Slide 20: Word Similarity Examples}
		\fontsize{6.5}{7.5}\selectfont
		\textbf{Similar Word Pairs:}
		\begin{tabular}{|c|c|c|}
			\hline
			\textbf{Word 1} & \textbf{Word 2} & \textbf{Similarity} \\
			\hline
			China & Japan & 0.84 \\
			King & Queen & 0.78 \\
			Gold & Silver & 0.95 \\
			Europe & Asia & 0.83 \\
			\hline
		\end{tabular}
		
		\vspace{0.1cm}
		\textbf{Dissimilar Word Pairs:}
		\begin{tabular}{|c|c|c|}
			\hline
			\textbf{Word 1} & \textbf{Word 2} & \textbf{Similarity} \\
			\hline
			Australia & Tokyo & 0.40 \\
			King & Man & 0.53 \\
			Gold & Dollar & 0.50 \\
			Banana & Rose & 0.30 \\
			\hline
		\end{tabular}
		
		\vspace{0.1cm}
		\textbf{Key Observation:}
		\begin{itemize}
			\item Similar categories = high similarity
			\item Different categories = low similarity
			\item Embeddings capture meaningful relationships
		\end{itemize}
		
		\textcolor{myblue}{\textbf{Key Insight:}} Embeddings distinguish similar and dissimilar words!
	\end{frame}
	
	\begin{frame}{Slide 21: Word Analogies}
		\fontsize{6.5}{7.5}\selectfont
		\textbf{Word Analogies:}
		\begin{itemize}
			\item "a is to b" as "c is to d"
			\item Captures relationships between word pairs
			\item Operations in embedding space
		\end{itemize}
		
		\vspace{0.1cm}
		\textbf{Examples:}
		\begin{tabular}{|c|c|c|}
			\hline
			\textbf{Operation} & \textbf{Closest Word} & \textbf{Similarity} \\
			\hline
			America - Washington + London & UK & 0.77 \\
			King - Man + Woman & Queen & 0.86 \\
			Gold - Dollar + Pound & Silver & 0.80 \\
			Europe - Spain + China & Asia & 0.79 \\
			\hline
		\end{tabular}
		
		\vspace{0.1cm}
		\textbf{How It Works:}
		\begin{enumerate}
			\item Embedding vectors capture relationships
			\item Vector arithmetic preserves relationships
			\item "king - man + woman equals queen"
		\end{enumerate}
		
		\textcolor{myblue}{\textbf{Key Insight:}} Vector arithmetic reveals semantic relationships!
	\end{frame}
	
	\begin{frame}{Slide 22: Visualizing Embeddings}
		\fontsize{6.5}{7.5}\selectfont
		\textbf{t-SNE Visualization:}
		\begin{itemize}
			\item t-SNE (t-distributed Stochastic Neighbor Embedding)
			\item Reduces embeddings to 2 dimensions
			\item Visualizes relationships between words
			\item Clusters show semantic groups
		\end{itemize}
		
		\vspace{0.1cm}
		\textbf{What We See:}
		\begin{enumerate}
			\item \textbf{Europe and Asia} - close together (continents)
			\item \textbf{Spain and China} - close together (countries)
			\item \textbf{Europe-Spain relationship} $\approx$ Asia-China relationship
		\end{enumerate}
		
		\vspace{0.1cm}
		\textbf{Key Insights from Visualization:}
		\begin{itemize}
			\item Related words cluster together
			\item Relationships are preserved in 2D
			\item Reveals structure in embeddings
		\end{itemize}
		
		\textcolor{myblue}{\textbf{Key Concept:}} t-SNE visualizes embedding relationships!
	\end{frame}
	
	\begin{frame}{Slide 23: Word Embedding Limitations}
		\fontsize{6.5}{7.5}\selectfont
		\textbf{Major Limitations:}
		\begin{enumerate}
			\item \textbf{No Sentence-Level Context}
			\begin{itemize}
				\item Same word has different meanings
				\item Example: "bat" (animal) vs "bat" (sports)
				\item Word2Vec gives same vector for both
			\end{itemize}
			\item \textbf{Fixed Context Window}
			\begin{itemize}
				\item Must specify window size a priori
				\item May miss distant relationships
				\item Limited by window size
			\end{itemize}
			\item \textbf{Ignore Word Order}
			\begin{itemize}
				\item Within the context window
				\item Loss of syntactic information
			\end{itemize}
		\end{enumerate}
		
		\textcolor{myblue}{\textbf{Example:}}
		\begin{itemize}
			\item "The bat flew out of the cave" vs "He picked up the bat to play baseball"
			\item Same embedding for both "bat" instances
			\item Context-dependent meaning is lost
		\end{itemize}
	\end{frame}
	
	\begin{frame}{Slide 24: Other Embedding Approaches}
		\fontsize{6.5}{7.5}\selectfont
		\textbf{Alternative Word Embeddings:}
		\begin{enumerate}
			\item \textbf{GloVe (Global Vectors)}
			\begin{itemize}
				\item Stanford University
				\item Combines matrix factorization with word2vec
				\item Uses global co-occurrence statistics
			\end{itemize}
			\item \textbf{Sentence Embeddings}
			\begin{itemize}
				\item SBert (Sentence BERT)
				\item Captures sentence-level meaning
				\item Better for sentence comparison
			\end{itemize}
			\item \textbf{Document Embeddings}
			\begin{itemize}
				\item Doc2Vec
				\item Creates document-level vectors
				\item Captures document meaning
			\end{itemize}
		\end{enumerate}
		
		\vspace{0.1cm}
		\textbf{Gemini Embeddings:}
		\begin{itemize}
			\item Advanced embedding model
			\item Part of Google's Gemini family
			\item Multimodal capabilities
		\end{itemize}
		
		\textcolor{myblue}{\textbf{Key Insight:}} Many embedding approaches exist for different needs!
	\end{frame}
	
	\begin{frame}{Slide 25: Embedding Applications}
		\fontsize{6.5}{7.5}\selectfont
		\textbf{Standard Input for NLP Systems:}
		\begin{enumerate}
			\item \textbf{Clustering}
			\begin{itemize}
				\item Group similar documents
				\item Topic identification
				\item Document categorization
			\end{itemize}
			\item \textbf{Sentiment Analysis}
			\begin{itemize}
				\item Customer feedback analysis
				\item Market sentiment detection
				\item Social media monitoring
			\end{itemize}
			\item \textbf{Recurrent Models}
			\begin{itemize}
				\item LSTM and RNN inputs
				\item Sequence processing
				\item Text generation
			\end{itemize}
			\item \textbf{Transformer Models}
			\begin{itemize}
				\item Pre-trained embeddings
				\item Transfer learning
				\item Contextual embeddings
			\end{itemize}
		\end{enumerate}
		
		\textcolor{myblue}{\textbf{Key Insight:}} Embeddings are foundation for modern NLP!
	\end{frame}
	
	% ============================================================
	% SECTION 3: RNNs AND LSTMs - SLIDES 26-35
	% ============================================================
	
	\begin{frame}{Slide 26: RNN Introduction}
		\fontsize{6.5}{7.5}\selectfont
		\textbf{What are RNNs?}
		\begin{itemize}
			\item Recurrent Neural Networks
			\item Developed for sequential data
			\item Handle text and speech
			\item Laid groundwork for GenAI
		\end{itemize}
		
		\vspace{0.1cm}
		\textbf{Key Features:}
		\begin{enumerate}
			\item \textbf{Memory Cell}
			\begin{itemize}
				\item Hidden cell stores information
				\item Value at time t depends on previous values
			\end{itemize}
			\item \textbf{Sequential Processing}
			\begin{itemize}
				\item Processes one word at a time
				\item Combines current inputs with memory
			\end{itemize}
			\item \textbf{Autoregressive Nature}
			\begin{itemize}
				\item Outputs become inputs for next step
				\item Captures dependencies
			\end{itemize}
		\end{enumerate}
		
		\textcolor{myblue}{\textbf{Key Insight:}} RNNs process sequential data with memory!
	\end{frame}
	
	\begin{frame}{Slide 27: RNN Architecture}
		\fontsize{6.5}{7.5}\selectfont
		\textbf{Basic RNN Structure:}
		\begin{itemize}
			\item Input at time t: $x_t$
			\item Hidden state at time t: $h_t$
			\item Output at time t: $y_t$
			\item $h_t = f(W_{xh}x_t + W_{hh}h_{t-1} + b_h)$
			\item $y_t = g(W_{hy}h_t + b_y)$
		\end{itemize}
		
		\vspace{0.1cm}
		\textbf{Diagram:}
		\begin{itemize}
			\item Input $\to$ Hidden Cell $\to$ Output
			\item Hidden cell feeds back to itself
			\item Previous hidden state influences current output
		\end{itemize}
		
		\vspace{0.1cm}
		\textbf{Advantages:}
		\begin{enumerate}
			\item Handles variable length sequences
			\item Captures order of words
			\item Learns dependencies
			\item Uses autoregressive structure
		\end{enumerate}
		
		\textcolor{myblue}{\textbf{Key Concept:}} RNNs combine current input with previous state!
	\end{frame}
	
	\begin{frame}{Slide 28: RNN vs Feedforward NN}
		\fontsize{6.5}{7.5}\selectfont
		\textbf{Feedforward Neural Network (FNN):}
		\begin{itemize}
			\item No memory
			\item Input features flow one direction
			\item Output depends only on current inputs
			\item Used for static data
		\end{itemize}
		
		\vspace{0.1cm}
		\textbf{Recurrent Neural Network (RNN):}
		\begin{itemize}
			\item Has memory (hidden state)
			\item Information can flow in loops
			\item Output depends on current inputs AND previous states
			\item Used for sequential data
		\end{itemize}
		
		\vspace{0.1cm}
		\textbf{Key Difference:}
		\begin{itemize}
			\item FNN: $y = f(x)$
			\item RNN: $y_t = f(x_t, h_{t-1})$
		\end{itemize}
		
		\textcolor{myblue}{\textbf{Remember:}} RNNs have memory; FNNs do not!
	\end{frame}
	
	\begin{frame}{Slide 29: Vanishing Gradient Problem in RNNs}
		\fontsize{6.5}{7.5}\selectfont
		\textbf{The Problem:}
		\begin{itemize}
			\item Gradients become very small during backpropagation
			\item Early layers receive tiny gradients
			\item Weights in early layers barely update
			\item Long-range dependencies are lost
		\end{itemize}
		
		\vspace{0.1cm}
		\textbf{Why It Happens:}
		\begin{enumerate}
			\item Gradients multiplied through many time steps
			\item Each multiplication by $< 1$
			\item Product approaches zero
			\item Repeated with each time step
		\end{enumerate}
		
		\vspace{0.1cm}
		\textbf{Example:}
		\begin{itemize}
			\item Sentence: "the analyst stopped using MGARCH because it is too hard to estimate"
			\item "MGARCH" is 7 words away from "estimate"
			\item RNN loses connection between them
		\end{itemize}
		
		\textcolor{myblue}{\textbf{Key Insight:}} Vanishing gradients limit RNN effectiveness!
	\end{frame}
	
	\begin{frame}{Slide 30: Long Short-Term Memory (LSTM)}
		\fontsize{6.5}{7.5}\selectfont
		\textbf{What is LSTM?}
		\begin{itemize}
			\item Special class of RNNs
			\item Developed to address vanishing gradient problem
			\item Uses layers/gates to regulate information flow
			\item Can capture long-range dependencies
		\end{itemize}
		
		\vspace{0.1cm}
		\textbf{Key Feature:}
		\begin{itemize}
			\item Learns what to remember
			\item Learns what to forget
			\item More sophisticated memory management
		\end{itemize}
		
		\vspace{0.1cm}
		\textbf{Advantages:}
		\begin{enumerate}
			\item Handles long-term dependencies
			\item More effective than basic RNNs
			\item Successful in many applications
			\item Used for time series, text, audio
		\end{enumerate}
		
		\textcolor{myblue}{\textbf{Key Insight:}} LSTMs solve the vanishing gradient problem!
	\end{frame}
	
	\begin{frame}{Slide 31: LSTM Architecture - The Gates}
		\fontsize{6.5}{7.5}\selectfont
		\textbf{Three Gates in LSTM:}
		\begin{enumerate}
			\item \textbf{Forget Gate}
			\begin{itemize}
				\item Determines what to discard from memory
				\item Forgets irrelevant information
				\item $f_t = \sigma(W_f \cdot [h_{t-1}, x_t] + b_f)$
			\end{itemize}
			\item \textbf{Input Gate}
			\begin{itemize}
				\item Determines what new information to store
				\item Updates memory with important information
				\item $i_t = \sigma(W_i \cdot [h_{t-1}, x_t] + b_i)$
			\end{itemize}
			\item \textbf{Output Gate}
			\begin{itemize}
				\item Determines what to output
				\item Uses relevant memory for output
				\item $o_t = \sigma(W_o \cdot [h_{t-1}, x_t] + b_o)$
			\end{itemize}
		\end{enumerate}
		
		\textcolor{myblue}{\textbf{Key Concept:}} Gates control information flow in LSTM!
	\end{frame}
	
	\begin{frame}{Slide 32: LSTM Cell Diagram}
		\fontsize{6.5}{7.5}\selectfont
		\textbf{LSTM Cell Components:}
		\begin{itemize}
			\item Input: $x_t$ (current input), $h_{t-1}$ (previous hidden state), $c_{t-1}$ (previous cell state)
			\item Output: $h_t$ (current hidden state), $c_t$ (current cell state)
		\end{itemize}
		
		\vspace{0.1cm}
		\textbf{Gate Functions:}
		\begin{enumerate}
			\item \textbf{Cell State Update:}
			\begin{itemize}
				\item $\tilde{c}_t = \tanh(W_c \cdot [h_{t-1}, x_t] + b_c)$
				\item Candidate values for new cell state
			\end{itemize}
			\item \textbf{Cell State:}
			\begin{itemize}
				\item $c_t = f_t \cdot c_{t-1} + i_t \cdot \tilde{c}_t$
				\item Combines forgetting and new information
			\end{itemize}
			\item \textbf{Hidden State:}
			\begin{itemize}
				\item $h_t = o_t \cdot \tanh(c_t)$
				\item Output based on filtered cell state
			\end{itemize}
		\end{enumerate}
		
		\textcolor{myblue}{\textbf{Key Insight:}} LSTM manages long-term memory effectively!
	\end{frame}
	
	\begin{frame}{Slide 33: LSTM vs RNN Comparison}
		\fontsize{6.5}{7.5}\selectfont
		\textbf{Basic RNN:}
		\begin{itemize}
			\item Simple structure
			\item Prone to vanishing gradients
			\item Struggles with long-range dependencies
			\item One activation function
		\end{itemize}
		
		\vspace{0.1cm}
		\textbf{LSTM:}
		\begin{itemize}
			\item Complex structure with gates
			\item Solves vanishing gradient
			\item Handles long-range dependencies
			\item Multiple activation functions
		\end{itemize}
		
		\vspace{0.1cm}
		\textbf{When to Use Each:}
		\begin{tabular}{|c|c|}
			\hline
			\textbf{RNN} & \textbf{LSTM} \\
			\hline
			Short sequences & Long sequences \\
			Simple patterns & Complex patterns \\
			Limited data & Large data \\
			Quick prototyping & Production models \\
			\hline
		\end{tabular}
		
		\textcolor{myblue}{\textbf{Key Insight:}} LSTM is more powerful but more complex!
	\end{frame}
	
	\begin{frame}{Slide 34: RNNs and LSTMs Applications}
		\fontsize{6.5}{7.5}\selectfont
		\textbf{Applications in Finance:}
		\begin{enumerate}
			\item \textbf{Time Series Analysis}
			\begin{itemize}
				\item Stock price prediction
				\item Volatility forecasting
				\item Economic indicator forecasting
			\end{itemize}
			\item \textbf{Text Processing}
			\begin{itemize}
				\item Document classification
				\item Sentiment analysis
				\item Language modeling
			\end{itemize}
			\item \textbf{Speech Recognition}
			\begin{itemize}
				\item Audio-to-text conversion
				\item Voice commands
				\item Speech synthesis
			\end{itemize}
		\end{enumerate}
		
		\vspace{0.1cm}
		\textbf{Why They Work:}
		\begin{itemize}
			\item Capture temporal dependencies
			\item Handle variable length sequences
			\item Learn from sequential patterns
		\end{itemize}
		
		\textcolor{myblue}{\textbf{Key Insight:}} RNNs/LSTMs excel at sequence-based tasks!
	\end{frame}
	
	\begin{frame}{Slide 35: RNN Limitations Summary}
		\fontsize{6.5}{7.5}\selectfont
		\textbf{Key Limitations:}
		\begin{enumerate}
			\item \textbf{Vanishing/Exploding Gradients}
			\begin{itemize}
				\item Limits long-range learning
				\item Requires careful design
			\end{itemize}
			\item \textbf{Sequential Processing}
			\begin{itemize}
				\item Can't parallelize
				\item Slower than transformers
			\end{itemize}
			\item \textbf{Memory Limitations}
			\begin{itemize}
				\item Forgets distant information
				\item Limited context window
			\end{itemize}
		\end{enumerate}
		
		\vspace{0.1cm}
		\textbf{The Solution:}
		\begin{itemize}
			\item Transformers address these limitations
			\item Self-attention mechanisms
			\item Parallel processing
			\item Better long-range dependencies
		\end{itemize}
		
		\textcolor{myblue}{\textbf{Key Insight:}} Transformers overcame RNN limitations!
	\end{frame}
	
	% ============================================================
	% SECTION 4: TRANSFORMERS - SLIDES 36-50
	% ============================================================
	
	\begin{frame}{Slide 36: Transformers Introduction}
		\fontsize{6.5}{7.5}\selectfont
		\textbf{What are Transformers?}
		\begin{itemize}
			\item Next major development in NLP
			\item Class of feedforward neural networks
			\item Based on "attention" mechanism
			\item Introduced in "Attention is All You Need" (2017)
		\end{itemize}
		
		\vspace{0.1cm}
		\textbf{Key Innovation:}
		\begin{itemize}
			\item Self-attention mechanism
			\item Multi-head attention
			\item Semantic understanding
			\item Contextual relationships
		\end{itemize}
		
		\vspace{0.1cm}
		\textbf{Why Transformers Matter:}
		\begin{enumerate}
			\item Superior to RNNs for NLP tasks
			\item Parallel processing
			\item Captures long-range dependencies
			\item Foundation for LLMs
		\end{enumerate}
		
		\textcolor{myblue}{\textbf{Key Insight:}} Transformers revolutionized NLP!
	\end{frame}
	
	\begin{frame}{Slide 37: Transformer Architecture Overview}
		\fontsize{6.5}{7.5}\selectfont
		\textbf{Two Main Components:}
		\begin{enumerate}
			\item \textbf{Encoder}
			\begin{itemize}
				\item Processes input text
				\item Converts to vectors
				\item "Understands" the text
			\end{itemize}
			\item \textbf{Decoder}
			\begin{itemize}
				\item Generates output text
				\item Based on input/prompt
				\item Creates responses
			\end{itemize}
		\end{enumerate}
		
		\vspace{0.1cm}
		\textbf{Evolution:}
		\begin{itemize}
			\item Some transformers specialize in encoding (BERT)
			\item Some focus on decoding (GPT)
			\item Some have both (BART)
		\end{itemize}
		
		\vspace{0.1cm}
		\textbf{Key Components:}
		\begin{itemize}
			\item Embedding layers
			\item Positional encoding
			\item Multi-head attention
			\item Feed-forward networks
			\item Layer normalization
		\end{itemize}
		
		\textcolor{myblue}{\textbf{Key Concept:}} Encoder processes; Decoder generates!
	\end{frame}
	
	\begin{frame}{Slide 38: Self-Attention Mechanism}
		\fontsize{6.5}{7.5}\selectfont
		\textbf{What is Self-Attention?}
		\begin{itemize}
			\item Calculates relationships between all words
			\item Creates semantic understanding
			\item Captures long-range dependencies
		\end{itemize}
		
		\vspace{0.1cm}
		\textbf{How It Works:}
		\begin{enumerate}
			\item \textbf{Query Vector}
			\begin{itemize}
				\item Current token being compared
				\item Asks "what should I pay attention to?"
			\end{itemize}
			\item \textbf{Key Vector}
			\begin{itemize}
				\item Other tokens being compared
				\item Response to the query
			\end{itemize}
			\item \textbf{Value Vector}
			\begin{itemize}
				\item Content of the token
				\item Used to create context vector
			\end{itemize}
		\end{enumerate}
		
		\textcolor{myblue}{\textbf{Key Insight:}} Attention determines what's important!
	\end{frame}
	
	\begin{frame}{Slide 39: Attention Calculation Steps}
		\fontsize{6.5}{7.5}\selectfont
		\textbf{Step 1: Score Calculation}
		\begin{itemize}
			\item Compare query and key vectors
			\item $score(q_i, k_j) = q_i \cdot k_j$
			\item Higher score = more relevant
		\end{itemize}
		
		\vspace{0.1cm}
		\textbf{Step 2: Softmax Normalization}
		\begin{itemize}
			\item $\alpha_{ij} = \text{softmax}(score(q_i, k_j))$
			\item Normalizes scores to probabilities
			\item Ensures gradients don't vanish
		\end{itemize}
		
		\vspace{0.1cm}
		\textbf{Step 3: Context Vector}
		\begin{itemize}
			\item $c_i = \sum_{j=1}^n \alpha_{ij} v_j$
			\item Weighted sum of value vectors
			\item Context-enriched embedding
		\end{itemize}
		
		\vspace{0.1cm}
		\textbf{Example:}
		\begin{itemize}
			\item "debt is senior to equity, so it is less risky"
			\item "it" has highest attention to "debt"
			\item Model learns "it" refers to "debt"
		\end{itemize}
		
		\textcolor{myblue}{\textbf{Key Insight:}} Attention = query-key similarity + weighted values!
	\end{frame}
	
	\begin{frame}{Slide 40: Multi-Head Attention}
		\fontsize{6.5}{7.5}\selectfont
		\textbf{What is Multi-Head Attention?}
		\begin{itemize}
			\item Multiple attention heads in parallel
			\item Each head focuses on different aspects
			\item Captures different relationships
		\end{itemize}
		
		\vspace{0.1cm}
		\textbf{How It Works:}
		\begin{enumerate}
			\item Partition input into sub-matrices
			\item Each head processes independently
			\item Combine results from all heads
		\end{enumerate}
		
		\vspace{0.1cm}
		\textbf{Example:}
		\begin{itemize}
			\item GPT basic version has 12 heads
			\item Each head captures different patterns
			\item Syntax, semantics, context, etc.
		\end{itemize}
		
		\vspace{0.1cm}
		\textbf{Benefits:}
		\begin{itemize}
			\item Richer representations
			\item Multiple relationship types
			\item Better understanding
		\end{itemize}
		
		\textcolor{myblue}{\textbf{Key Insight:}} Multi-head captures diverse relationships!
	\end{frame}
	
	\begin{frame}{Slide 41: Positional Embeddings}
		\fontsize{6.5}{7.5}\selectfont
		\textbf{Why Positional Information Matters:}
		\begin{itemize}
			\item Word order is crucial for meaning
			\item "dog bites man" $\neq$ "man bites dog"
			\item Need to capture position in sequence
		\end{itemize}
		
		\vspace{0.1cm}
		\textbf{Two Approaches:}
		\begin{enumerate}
			\item \textbf{Absolute Positional Embeddings}
			\begin{itemize}
				\item Sine and cosine functions
				\item $PE_{p,2i} = \sin(p/10000^{2i/d})$
				\item $PE_{p,2i+1} = \cos(p/10000^{2i/d})$
			\end{itemize}
			\item \textbf{Relative Positional Embeddings}
			\begin{itemize}
				\item Adds position number to attention score
				\item Focuses on relative distances
			\end{itemize}
		\end{enumerate}
		
		\vspace{0.1cm}
		\textbf{Parametric Approach:}
		\begin{itemize}
			\item Initialize with random numbers
			\item Learn during training
			\item Flexible and adaptive
		\end{itemize}
		
		\textcolor{myblue}{\textbf{Key Insight:}} Positional embeddings preserve word order!
	\end{frame}
	
	\begin{frame}{Slide 42: Transformer Advantages}
		\fontsize{6.5}{7.5}\selectfont
		\textbf{Advantages over RNNs:}
		\begin{enumerate}
			\item \textbf{Parallel Processing}
			\begin{itemize}
				\item Processes entire sequence at once
				\item Much faster training
				\item Better hardware utilization
			\end{itemize}
			\item \textbf{Long-Range Dependencies}
			\begin{itemize}
				\item Direct connections between words
				\item No information loss over distance
				\item Better semantic understanding
			\end{itemize}
			\item \textbf{Contextualized Embeddings}
			\begin{itemize}
				\item Different vectors for different contexts
				\item Same word, different meaning
				\item Handles polysemy effectively
			\end{itemize}
			\item \textbf{Universal World Representations}
			\begin{itemize}
				\item General knowledge
				\item Multiple tasks
				\item Transfer learning
			\end{itemize}
		\end{enumerate}
		
		\textcolor{myblue}{\textbf{Key Insight:}} Transformers are superior for most NLP tasks!
	\end{frame}
	
	\begin{frame}{Slide 43: Transformer Disadvantages}
		\fontsize{6.5}{7.5}\selectfont
		\textbf{Key Disadvantages:}
		\begin{enumerate}
			\item \textbf{Computational Intensity}
			\begin{itemize}
				\item More computationally intensive than RNNs
				\item Requires powerful hardware
				\item Higher energy consumption
				\item Longer training times
			\end{itemize}
			\item \textbf{Interpretability Challenges}
			\begin{itemize}
				\item Complex, multi-layered structure
				\item Harder to understand decisions
				\item Difficult to identify error sources
				\item Opaque reasoning
			\end{itemize}
			\item \textbf{Data Requirements}
			\begin{itemize}
				\item Need huge training datasets
				\item Billions of parameters
				\item Extensive computation
			\end{itemize}
		\end{enumerate}
		
		\textcolor{myblue}{\textbf{Key Insight:}} Transformers are powerful but resource-intensive!
	\end{frame}
	
	\begin{frame}{Slide 44: Attention Types}
		\fontsize{6.5}{7.5}\selectfont
		\textbf{Self-Attention:}
		\begin{itemize}
			\item Words analyze other words in same sequence
			\item Captures internal relationships
			\item Example: "it" $\to$ "debt"
		\end{itemize}
		
		\vspace{0.1cm}
		\textbf{Cross-Attention:}
		\begin{itemize}
			\item Words analyze words in another sequence
			\item Used in encoder-decoder models
			\item Example: Translation tasks
		\end{itemize}
		
		\vspace{0.1cm}
		\textbf{Masked Attention:}
		\begin{itemize}
			\item Blocks future tokens in decoder
			\item Prevents cheating
			\item Used for next-token prediction
		\end{itemize}
		
		\vspace{0.1cm}
		\textbf{Historical Attention Methods:}
		\begin{enumerate}
			\item Additive Attention (Bahdanau)
			\item Dot-Product Attention (Luong)
			\item Scaled Dot-Product Attention (Vaswani)
		\end{enumerate}
		
		\textcolor{myblue}{\textbf{Key Concept:}} Different attention types for different tasks!
	\end{frame}
	
	\begin{frame}{Slide 45: Transformer Layers Architecture}
		\fontsize{6.5}{7.5}\selectfont
		\textbf{Encoder Layer:}
		\begin{enumerate}
			\item \textbf{Multi-Head Self-Attention}
			\begin{itemize}
				\item Captures relationships between input tokens
				\item Produces context vectors
			\end{itemize}
			\item \textbf{Add \& Normalize}
			\begin{itemize}
				\item Residual connection
				\item Prevents vanishing gradients
			\end{itemize}
			\item \textbf{Feed-Forward Network}
			\begin{itemize}
				\item Processes context vectors
				\item Nonlinear transformations
			\end{itemize}
			\item \textbf{Add \& Normalize}
			\begin{itemize}
				\item Another residual connection
				\item Layer normalization
			\end{itemize}
		\end{enumerate}
		
		\vspace{0.1cm}
		\textbf{Decoder Layer:}
		\begin{itemize}
			\item Similar structure plus masked attention
			\item Prevents seeing future tokens
			\item Autoregressive generation
		\end{itemize}
		
		\textcolor{myblue}{\textbf{Key Insight:}} Multiple layers with residual connections!
	\end{frame}
	
	\begin{frame}{Slide 46: Transformer vs RNN Comparison}
		\fontsize{6.5}{7.5}\selectfont
		\textbf{Processing Style:}
		\begin{itemize}
			\item \textbf{RNN:} Sequential (word by word)
			\item \textbf{Transformer:} Parallel (all words at once)
		\end{itemize}
		
		\vspace{0.1cm}
		\textbf{Memory/Context:}
		\begin{itemize}
			\item \textbf{RNN:} Limited by hidden state
			\item \textbf{Transformer:} Direct connections to all words
		\end{itemize}
		
		\vspace{0.1cm}
		\textbf{Training Speed:}
		\begin{itemize}
			\item \textbf{RNN:} Slow (sequential)
			\item \textbf{Transformer:} Fast (parallel)
		\end{itemize}
		
		\vspace{0.1cm}
		\textbf{Long-Range Dependencies:}
		\begin{itemize}
			\item \textbf{RNN:} Struggles (vanishing gradients)
			\item \textbf{Transformer:} Excels (attention)
		\end{itemize}
		
		\vspace{0.1cm}
		\textbf{Interpretability:}
		\begin{itemize}
			\item \textbf{RNN:} More interpretable
			\item \textbf{Transformer:} Less interpretable
		\end{itemize}
		
		\textcolor{myblue}{\textbf{Key Insight:}} Transformers win on performance; RNNs on interpretability!
	\end{frame}
	
	\begin{frame}{Slide 47: BERT Architecture}
		\fontsize{6.5}{7.5}\selectfont
		\textbf{What is BERT?}
		\begin{itemize}
			\item Bidirectional Encoder Representations from Transformers
			\item Introduced by Google in 2018
			\item Encoder-only architecture
			\item ~100 million parameters
		\end{itemize}
		
		\vspace{0.1cm}
		\textbf{Key Features:}
		\begin{enumerate}
			\item \textbf{Bidirectional}
			\begin{itemize}
				\item Examines words before and after
				\item Better context understanding
				\item More accurate representations
			\end{itemize}
			\item \textbf{Training Methods}
			\begin{itemize}
				\item Masked Language Modeling (15\% masked)
				\item Next Sentence Prediction
			\end{itemize}
			\item \textbf{Applications}
			\begin{itemize}
				\item Classification
				\item Sentiment analysis
				\item Question answering
			\end{itemize}
		\end{enumerate}
		
		\textcolor{myblue}{\textbf{Key Insight:}} BERT is encoder-only and bidirectional!
	\end{frame}
	
	\begin{frame}{Slide 48: BERT Training}
		\fontsize{6.5}{7.5}\selectfont
		\textbf{Stage 1: Pre-training}
		\begin{itemize}
			\item Trained on large unlabeled corpus
			\item Books and Wikipedia pages
			\item Learned general language understanding
		\end{itemize}
		
		\vspace{0.1cm}
		\textbf{Stage 2: Self-Supervision}
		\begin{enumerate}
			\item \textbf{Masked Language Modeling}
			\begin{itemize}
				\item 15\% of words randomly masked
				\item Model predicts hidden words
				\item Learns bidirectional context
			\end{itemize}
			\item \textbf{Next Sentence Prediction}
			\begin{itemize}
				\item Sentence pairs from documents
				\item Predict if sentences are sequential
				\item Understands relationships
			\end{itemize}
		\end{enumerate}
		
		\vspace{0.1cm}
		\textbf{Result:}
		\begin{itemize}
			\item Deep language understanding
			\item Transferable knowledge
			\item Strong performance on NLP tasks
		\end{itemize}
		
		\textcolor{myblue}{\textbf{Key Insight:}} BERT learns from masked words and sentence pairs!
	\end{frame}
	
	\begin{frame}{Slide 49: BERT Variants}
		\fontsize{6.5}{7.5}\selectfont
		\textbf{RoBERTa:}
		\begin{itemize}
			\item Robustly Optimized BERT Pre-training Approach
			\item Different masking patterns each cycle
			\item Larger training dataset
			\item More training cycles
			\item Multi-language web crawled data
		\end{itemize}
		
		\vspace{0.1cm}
		\textbf{FinBERT:}
		\begin{itemize}
			\item Financial text specialization
			\item Trained on Reuters financial documents
			\item Corporate statements
			\item Analyst reports
			\item Optimal for financial sentiment analysis
		\end{itemize}
		
		\vspace{0.1cm}
		\textbf{Key Difference:}
		\begin{itemize}
			\item RoBERTa: General improvement
			\item FinBERT: Domain specialization
		\end{itemize}
		
		\textcolor{myblue}{\textbf{Key Insight:}} Specialized BERT variants for specific tasks!
	\end{frame}
	
	\begin{frame}{Slide 50: BART Architecture}
		\fontsize{6.5}{7.5}\selectfont
		\textbf{What is BART?}
		\begin{itemize}
			\item Bidirectional and Auto-Regressive Transformer
			\item Developed by Facebook (Meta) in 2019
			\item Encoder-Decoder architecture
			\item More versatile than BERT
		\end{itemize}
		
		\vspace{0.1cm}
		\textbf{Key Features:}
		\begin{enumerate}
			\item \textbf{Bidirectional Encoder}
			\begin{itemize}
				\item Similar to BERT
				\item Understands input
			\end{itemize}
			\item \textbf{Auto-Regressive Decoder}
			\begin{itemize}
				\item Similar to GPT
				\item Generates text
			\end{itemize}
			\item \textbf{Training Method}
			\begin{itemize}
				\item Distorts blocks of text
				\item Shuffles, removes, replaces words
				\item Learns to generate original forms
			\end{itemize}
		\end{enumerate}
		
		\textcolor{myblue}{\textbf{Key Insight:}} BART combines BERT and GPT capabilities!
	\end{frame}
	
	% ============================================================
	% SECTION 5: LARGE LANGUAGE MODELS - SLIDES 51-70
	% ============================================================
	
	\begin{frame}{Slide 51: LLM Definition}
		\fontsize{6.5}{7.5}\selectfont
		\textbf{What are Large Language Models?}
		\begin{itemize}
			\item Predominantly based on transformers
			\item Pre-trained on vast textual datasets
			\item Large neural networks with billions of parameters
			\item Trained on very large corpora
		\end{itemize}
		
		\vspace{0.1cm}
		\textbf{Key Characteristics:}
		\begin{enumerate}
			\item \textbf{Size}
			\begin{itemize}
				\item Billions to trillions of parameters
				\item Huge computational requirements
			\end{itemize}
			\item \textbf{Training}
			\begin{itemize}
				\item Trained only once
				\item On massive text corpus
			\end{itemize}
			\item \textbf{Storage}
			\begin{itemize}
				\item Cloud-based (too large for local)
				\item Remote access
			\end{itemize}
		\end{enumerate}
		
		\textcolor{myblue}{\textbf{Key Insight:}} LLMs are huge transformer-based models!
	\end{frame}
	
	\begin{frame}{Slide 52: LLM Types}
		\fontsize{6.5}{7.5}\selectfont
		\textbf{Autoencoding LLMs:}
		\begin{itemize}
			\item Encode sentences by masking input
			\item Predict masked words
			\item Example: BERT
			\item Good for understanding
		\end{itemize}
		
		\vspace{0.1cm}
		\textbf{Autoregressive LLMs:}
		\begin{itemize}
			\item Predicts next word/token
			\item Generates text
			\item Example: GPT family
			\item Good for generation
		\end{itemize}
		
		\vspace{0.1cm}
		\textbf{Combination LLMs:}
		\begin{itemize}
			\item Both encoder and decoder
			\item Most versatile
			\item Example: BART
			\item Good for many tasks
		\end{itemize}
		
		\textcolor{myblue}{\textbf{Key Insight:}} Different LLM types for different purposes!
	\end{frame}
	
	\begin{frame}{Slide 53: LLM Training Process}
		\fontsize{6.5}{7.5}\selectfont
		\textbf{Step 1: Data Preparation}
		\begin{itemize}
			\item Clean training data (corpus)
			\item Filter and remove duplicates
			\item Remove noisy data and punctuation
			\item Resolve or remove ambiguous data
			\item Balance class distributions
		\end{itemize}
		
		\vspace{0.1cm}
		\textbf{Step 2: Tokenization}
		\begin{itemize}
			\item Separate data into small parts (tokens)
			\item Enable efficient processing
			\item Store positional information
		\end{itemize}
		
		\vspace{0.1cm}
		\textbf{Step 3: Pre-training}
		\begin{itemize}
			\item Train on unlabeled text
			\item Understand word relationships
			\item Learn language patterns
		\end{itemize}
		
		\textcolor{myblue}{\textbf{Key Insight:}} Training requires massive data and compute!
	\end{frame}
	
	\begin{frame}{Slide 54: LLM Training Process (continued)}
		\fontsize{6.5}{7.5}\selectfont
		\textbf{Step 4: Fine-tuning}
		\begin{itemize}
			\item Transfer knowledge to specific tasks
			\item Use labeled data
			\item Low-shot task transfer
		\end{itemize}
		
		\vspace{0.1cm}
		\textbf{Step 5: Alignment}
		\begin{itemize}
			\item Conform to human preferences
			\item Reinforce with human feedback (RLHF)
			\item Based on user feedback
		\end{itemize}
		
		\vspace{0.1cm}
		\textbf{Step 6: Decoding}
		\begin{itemize}
			\item Generate output based on input
			\item Different strategies
			\item Quality optimization
		\end{itemize}
		
		\vspace{0.1cm}
		\textbf{Step 7: Production Optimization}
		\begin{itemize}
			\item Turn off lower-level layers
			\item Drop unimportant weights
			\item Reduce weight precision
		\end{itemize}
		
		\textcolor{myblue}{\textbf{Key Insight:}} Multiple stages from training to production!
	\end{frame}
	
	\begin{frame}{Slide 55: Cloud-based LLMs}
		\fontsize{6.5}{7.5}\selectfont
		\textbf{Why Cloud-based?}
		\begin{itemize}
			\item Corpus is huge
			\item Weight matrices are massive
			\item Training only done once
			\item Infeasible to store locally
		\end{itemize}
		
		\vspace{0.1cm}
		\textbf{Access Methods:}
		\begin{enumerate}
			\item API access
			\begin{itemize}
				\item Web-based interfaces
				\item Programmatic access
			\end{itemize}
			\item Remote storage
			\begin{itemize}
				\item Proprietary to developers
				\item Regular updates
			\end{itemize}
		\end{enumerate}
		
		\vspace{0.1cm}
		\textbf{Popular Cloud LLMs:}
		\begin{itemize}
			\item OpenAI GPT series
			\item Google Gemini
			\item Meta Llama
			\item Anthropic Claude
		\end{itemize}
		
		\textcolor{myblue}{\textbf{Key Insight:}} LLMs reside in cloud due to size and proprietary nature!
	\end{frame}
	
	\begin{frame}{Slide 56: GPT Family Overview}
		\fontsize{6.5}{7.5}\selectfont
		\textbf{GPT-1 (2018):}
		\begin{itemize}
			\item ~100 million parameters
			\item Pre-trained on unpublished books
			\item Fine-tuned for specific tasks		\end{itemize}
		
		\vspace{0.1cm}
		\textbf{GPT-2 (2019):}
		\begin{itemize}
			\item 124 million (base) to 774 million (large)
			\item No-shot fine-tuning possible
			\item Larger, more varied corpus
		\end{itemize}
		
		\vspace{0.1cm}
		\textbf{GPT-3 (2020):}
		\begin{itemize}
			\item 175 billion parameters
			\item Vast web text, books, Wikipedia
			\item Generative capacity
			\item Synthetic news creation
		\end{itemize}
		
		\vspace{0.1cm}
		\textbf{GPT-4 (2023):}
		\begin{itemize}
			\item Advanced capabilities
			\item Multimodal
			\item Improved reasoning
		\end{itemize}
		
		\textcolor{myblue}{\textbf{Key Insight:}} Each generation significantly improved!
	\end{frame}
	
	\begin{frame}{Slide 57: GPT Evolution}
		\fontsize{6.5}{7.5}\selectfont
		\textbf{Parameter Growth:}
		\begin{itemize}
			\item GPT-1: 100 million
			\item GPT-2: 124-774 million
			\item GPT-3: 175 billion
			\item GPT-4: Estimated > 1 trillion
		\end{itemize}
		
		\vspace{0.1cm}
		\textbf{Capability Growth:}
		\begin{enumerate}
			\item \textbf{Better Low-Shot Transfer}
			\begin{itemize}
				\item Fewer examples needed
				\item Improved generalization
			\end{itemize}
			\item \textbf{No-Shot Fine-Tuning}
			\begin{itemize}
				\item No examples needed
				\item Adapts to new situations
			\end{itemize}
			\item \textbf{Generative Capacity}
			\begin{itemize}
				\item Creates new text
				\item Human-like quality
			\end{itemize}
		\end{enumerate}
		
		\textcolor{myblue}{\textbf{Key Insight:}} Larger models = better performance!
	\end{frame}
	
	\begin{frame}{Slide 58: Chatbots}
		\fontsize{6.5}{7.5}\selectfont
		\textbf{What are Chatbots?}
		\begin{itemize}
			\item Software programs with conversational interface
			\item Interact with human users
			\item Web-based or phone-based
		\end{itemize}
		
		\vspace{0.1cm}
		\textbf{Types of Chatbots:}
		\begin{enumerate}
			\item \textbf{FAQ-Based}
			\begin{itemize}
				\item Standard responses
				\item Limited conversational ability
				\item Simple information retrieval
			\end{itemize}
			\item \textbf{Flow-Based}
			\begin{itemize}
				\item Interactive conversations
				\item Follow-up questions
				\item Multi-step interactions
			\end{itemize}
			\item \textbf{LLM-Powered}
			\begin{itemize}
				\item Deep understanding
				\item Informed answers
				\item Human-like responses
			\end{itemize}
		\end{enumerate}
		
		\textcolor{myblue}{\textbf{Key Insight:}} LLMs significantly improve chatbot quality!
	\end{frame}
	
	\begin{frame}{Slide 59: Prompt Engineering}
		\fontsize{6.5}{7.5}\selectfont
		\textbf{What is Prompt Engineering?}
		\begin{itemize}
			\item Art and science of designing prompts
			\item Effective use of LLMs
			\item Iterative process
			\item Getting best results
		\end{itemize}
		
		\vspace{0.1cm}
		\textbf{Guidelines for Good Prompts:}
		\begin{enumerate}
			\item Direct and concise
			\item Provide context and reference material
			\item Split tasks into smaller parts
			\item Provide sufficient information
			\item Be flexible for different models
			\item Allow model time to find results
		\end{enumerate}
		
		\vspace{0.1cm}
		\textbf{Best Practice:}
		\begin{itemize}
			\item Use few-shot learning
			\item Ask for structured output
			\item Test different personas
			\item Iterate and improve
		\end{itemize}
		
		\textcolor{myblue}{\textbf{Key Insight:}} Good prompts = good responses!
	\end{frame}
	
	\begin{frame}{Slide 60: Zero, One, and Few-Shot Learning}
		\fontsize{6.5}{7.5}\selectfont
		\textbf{Zero-Shot Prompting:}
		\begin{itemize}
			\item No examples provided
			\item Most common type
			\item User inputs question directly
			\item Relies on model's training
		\end{itemize}
		\textbf{Example:} "What is the capital of France?"
		
		\vspace{0.1cm}
		\textbf{One-Shot Prompting:}
		\begin{itemize}
			\item One example provided
			\item Shows expected format
			\item Guides model response
		\end{itemize}
		\textbf{Example:} "Question: What is 2+2? Answer: 4. Question: What is 3+3?"
		
		\vspace{0.1cm}
		\textbf{Few-Shot Prompting:}
		\begin{itemize}
			\item Multiple examples provided
			\item Better context
			\item More accurate responses
		\end{itemize}
		\textbf{Example:} Several Q\&A pairs before the target question
		
		\textcolor{myblue}{\textbf{Key Insight:}} More examples = better responses (usually)!
	\end{frame}
	
	\begin{frame}{Slide 61: System Prompting}
		\fontsize{6.5}{7.5}\selectfont
		\textbf{System Prompting:}
		\begin{itemize}
			\item Defines overall purpose
			\item Sets limits and constraints
			\item Specifies model role
			\item Provides context
		\end{itemize}
		
		\vspace{0.1cm}
		\textbf{Role Prompting:}
		\begin{itemize}
			\item Specifies persona for LLM
			\item "You are a financial analyst..."
			\item "You are an expert in risk management..."
			\item Shapes response style and content
		\end{itemize}
		
		\vspace{0.1cm}
		\textbf{Contextual Prompting:}
		\begin{itemize}
			\item Provides background information
			\item Sets the scene
			\item Informs the response
		\end{itemize}
		
		\textbf{Example:}
		"You are a risk manager at a large bank. Analyze the following credit risk scenario..."
		
		\textcolor{myblue}{\textbf{Key Insight:}} System prompts shape LLM behavior!
	\end{frame}
	
	\begin{frame}{Slide 62: Chain of Thought (CoT) Prompting}
		\fontsize{6.5}{7.5}\selectfont
		\textbf{What is CoT Prompting?}
		\begin{itemize}
			\item Prompting with intermediate reasoning steps
			\item Solves complex reasoning tasks
			\item Improves multi-step reasoning
			\item Better than direct prompting
		\end{itemize}
		
		\vspace{0.1cm}
		\textbf{Example - Standard Prompting:}
		Q: "A company sells a product for \$50 each. Fixed costs are \$10,000, variable costs are \$30 per product. How many units to break even?"
		A: "10,000" (incorrect!)
		
		\vspace{0.1cm}
		\textbf{Example - CoT Prompting:}
		"Let's think step by step..."
		The model then shows: Total Cost = \$10,000 + \$30x, Total Revenue = \$50x, Solve: \$50x = \$10,000 + \$30x, \$20x = \$10,000, x = 500 units
		
		\textcolor{myblue}{\textbf{Key Insight:}} Step-by-step reasoning yields correct answers!
	\end{frame}
	
	\begin{frame}{Slide 63: CoT Variations}
		\fontsize{6.5}{7.5}\selectfont
		\textbf{Zero-Shot CoT:}
		\begin{itemize}
			\item Model generates reasoning chain
			\item No examples shown
			\item Simple but effective
			\item "Let's think step by step"
		\end{itemize}
		
		\vspace{0.1cm}
		\textbf{Few-Shot CoT:}
		\begin{itemize}
			\item Examples with reasoning shown
			\item Guides model approach
			\item Better for complex tasks
		\end{itemize}
		\textbf{Example:}
		"Q1: If there are 3 red and 4 blue balls, probability red? A1: Total balls = 7, red = 3, probability = 3/7"
		
		\vspace{0.1cm}
		\textbf{Auto-CoT:}
		\begin{itemize}
			\item Automates example creation
			\item Clustering of questions
			\item Representative selection
			\item Generated reasoning chains
		\end{itemize}
		
		\textcolor{myblue}{\textbf{Key Insight:}} CoT works best for large models (>10B parameters)!
	\end{frame}
	
	\begin{frame}{Slide 64: Self-Consistency and Tree of Thought}
		\fontsize{6.5}{7.5}\selectfont
		\textbf{Self-Consistency:}
		\begin{itemize}
			\item Generate multiple responses
			\item Choose most consistent response
			\item Improves reliability
			\item Reduces randomness
		\end{itemize}
		
		\vspace{0.1cm}
		\textbf{How It Works:}
		\begin{enumerate}
			\item Generate several answers
			\item Compare responses
			\item Select most common/consistent
			\item Better than single response
		\end{enumerate}
		
		\vspace{0.1cm}
		\textbf{Tree of Thought (ToT):}
		\begin{itemize}
			\item Generalization of CoT
			\item Multiple reasoning threads
			\item Explore different paths
			\item Most consistent answer selected
		\end{itemize}
		
		\vspace{0.1cm}
		\textbf{Example:}
		\begin{itemize}
			\item CoT: One reasoning path
			\item ToT: Multiple paths explored
			\item Better for complex problems
		\end{itemize}
		
		\textcolor{myblue}{\textbf{Key Insight:}} Multiple answers = more reliable results!
	\end{frame}
	
	\begin{frame}{Slide 65: Context Length}
		\fontsize{6.5}{7.5}\selectfont
		\textbf{What is Context Length?}
		\begin{itemize}
			\item Maximum tokens in input/output sequence
			\item Also called context window
			\item Limits model's working memory
			\item Important for long documents
		\end{itemize}
		
		\vspace{0.1cm}
		\textbf{Evolution of Context Length:}
		\begin{tabular}{|c|c|c|}
			\hline
			\textbf{Model} & \textbf{Tokens} & \textbf{Approx. Pages} \\
			\hline
			GPT-1 & 512 & 1.5 \\
			GPT-4 & 32k & 95 \\
			Gemini 2.5 Pro & 1048k & 1,500 \\
			Llama 4 & 10 million & Over 10,000 \\
			\hline
		\end{tabular}
		
		\vspace{0.1cm}
		\textbf{Important Notes:}
		\begin{itemize}
			\item Effective context length often shorter
			\item Performance drops after effective range
			\item Very long inputs not ideal for real-world
		\end{itemize}
		
		\textcolor{myblue}{\textbf{Key Insight:}} Context length limits how much text model can process!
	\end{frame}
	
	\begin{frame}{Slide 66: Statelessness}
		\fontsize{6.5}{7.5}\selectfont
		\textbf{What is Statelessness?}
		\begin{itemize}
			\item LLMs don't retain information between prompts
			\item Each prompt processed independently
			\item No memory of previous conversations
			\item By design
		\end{itemize}
		
		\vspace{0.1cm}
		\textbf{Implications:}
		\begin{enumerate}
			\item \textbf{No Memory}
			\begin{itemize}
				\item Cannot recall previous questions
				\item Each interaction starts fresh
			\end{itemize}
			\item \textbf{Context Management}
			\begin{itemize}
				\item Must provide full context each time
				\item Feed previous responses to maintain context
			\end{itemize}
			\item \textbf{User Responsibility}
			\begin{itemize}
				\item Must manage conversation flow
				\item Provide relevant history
			\end{itemize}
		\end{enumerate}
		
		\vspace{0.1cm}
		\textbf{Solution:}
		\begin{itemize}
			\item Users feed previous responses to LLM
			\item Make it aware of conversation context
			\item Maintain continuity
		\end{itemize}
		
		\textcolor{myblue}{\textbf{Key Insight:}} LLMs don't remember past conversations!
	\end{frame}
	
	\begin{frame}{Slide 67: Temperature Parameter}
		\fontsize{6.5}{7.5}\selectfont
		\textbf{What is Temperature?}
		\begin{itemize}
			\item Controls randomness/creativity
			\item Shapes probability distribution
			\item Affects token selection
			\item Range varies by model (e.g., 0-2 for OpenAI)
		\end{itemize}
		
		\vspace{0.1cm}
		\textbf{Effects:}
		\begin{tabular}{|c|c|c|}
			\hline
			\textbf{Temperature} & \textbf{Effect} & \textbf{Use Case} \\
			\hline
			Below 1 & Narrowed distribution & Predictable, deterministic \\
			1 & Unaltered distribution & Balanced \\
			Above 1 & Widened distribution & Creative, varied \\
			\hline
		\end{tabular}
		
		\vspace{0.1cm}
		\textbf{Examples:}
		\begin{itemize}
			\item Temperature = 0: Always picks most likely token
			\item Temperature = 0.5: Some variety, fairly predictable
			\item Temperature = 1.0: Balanced creativity
			\item Temperature = 2.0: Very creative, may be nonsensical
		\end{itemize}
		
		\textcolor{myblue}{\textbf{Key Insight:}} Lower temperature = more predictable; higher = more creative!
	\end{frame}
	
	\begin{frame}{Slide 68: Top-K and Top-P Sampling}
		\fontsize{6.5}{7.5}\selectfont
		\textbf{Top-K Sampling:}
		\begin{itemize}
			\item Limits to K most probable tokens
			\item Restricts choices
			\item Low K = less creativity
			\item High K = more variety
		\end{itemize}
		\textbf{Example:} K=3 means only top 3 tokens considered
		
		\vspace{0.1cm}
		\textbf{Top-P Sampling:}
		\begin{itemize}
			\item Selects smallest set with cumulative probability $\ge$ P
			\item Dynamic selection
			\item Low P = less creativity
			\item High P = more variety
		\end{itemize}
		\textbf{Example:} P=0.85 includes tokens until probability sum reaches 85\%
		
		\vspace{0.1cm}
		\textbf{Example Tokens:}
		\begin{tabular}{|c|c|c|c|}
			\hline
			\textbf{Token} & \textbf{Prob} & \textbf{K=3} & \textbf{P=0.85} \\
			\hline
			strong & 0.35 & $\checkmark$ & $\checkmark$ \\
			positive & 0.25 & $\checkmark$ & $\checkmark$ \\
			mixed & 0.15 & $\checkmark$ & $\checkmark$ \\
			weak & 0.10 & - & $\checkmark$ \\
			\hline
		\end{tabular}
		
		\textcolor{myblue}{\textbf{Key Insight:}} Top-K/P control which tokens can be selected!
	\end{frame}
	
	\begin{frame}{Slide 69: Temperature, Top-K, Top-P Combinations}
		\fontsize{6.5}{7.5}\selectfont
		\textbf{Low Temperature + High K/P:}
		\begin{itemize}
			\item More predictable responses
			\item Some variety in tokens
			\item Good for consistent output
			\item Recommended for most use cases
		\end{itemize}
		
		\vspace{0.1cm}
		\textbf{High Temperature + Low K/P:}
		\begin{itemize}
			\item Less predictable
			\item Less variety in tokens
			\item More creative within limited set
			\item May be less coherent
		\end{itemize}
		
		\vspace{0.1cm}
		\textbf{Other Parameters:}
		\begin{itemize}
			\item Model version selection
			\item Penalty for repeated tokens
			\item Frequency penalty
			\item Presence penalty
		\end{itemize}
		
		\vspace{0.1cm}
		\textbf{Best Practice:}
		\begin{itemize}
			\item Start with defaults
			\item Adjust based on needs
			\item Test different combinations
		\end{itemize}
		
		\textcolor{myblue}{\textbf{Key Insight:}} Parameters work together to control output!
	\end{frame}
	
	\begin{frame}{Slide 70: Agentic AI}
		\fontsize{6.5}{7.5}\selectfont
		\textbf{What is Agentic AI?}
		\begin{itemize}
			\item Next frontier in AI
			\item AI agents that take inputs and perform actions autonomously
			\item Operate using LLM outputs
			\item Dynamic decision-making
		\end{itemize}
		
		\vspace{0.1cm}
		\textbf{Key Features:}
		\begin{enumerate}
			\item \textbf{Autonomous Action}
			\begin{itemize}
				\item Takes actions without human input
				\item Schedules meetings
				\item Sends invites
				\item Reschedules conflicts
			\end{itemize}
			\item \textbf{Environment Awareness}
			\begin{itemize}
				\item Observes environment
				\item Takes decisions
				\item Calls external resources
			\end{itemize}
			\item \textbf{Learning Capability}
			\begin{itemize}
				\item Learns on the job
				\item Adjusts dynamically
			\end{itemize}
		\end{enumerate}
		
		\textcolor{myblue}{\textbf{Key Insight:}} Agentic AI acts; Generative AI creates!
	\end{frame}
	
	% ============================================================
	% SECTION 6: GENAI PERFORMANCE ANALYSIS - SLIDES 71-85
	% ============================================================
	
	\begin{frame}{Slide 71: GenAI Evaluation Challenges}
		\fontsize{6.5}{7.5}\selectfont
		\textbf{Why GenAI Evaluation is Complex:}
		\begin{enumerate}
			\item \textbf{Open-Ended Output}
			\begin{itemize}
				\item Unstructured responses
				\item No single correct answer
				\item Many valid variations
			\end{itemize}
			\item \textbf{Factual Accuracy}
			\begin{itemize}
				\item Hallucination risk
				\item May sound plausible but be wrong
				\item Need verification
			\end{itemize}
			\item \textbf{Bias and Fairness}
			\begin{itemize}
				\item Potential for unethical output
				\item Bias in training data
				\item Fairness concerns
			\end{itemize}
			\item \textbf{Creativity vs Accuracy}
			\begin{itemize}
				\item Balancing creativity with correctness
				\item Task-dependent requirements
			\end{itemize}
		\end{enumerate}
		
		\textcolor{myblue}{\textbf{Key Insight:}} GenAI evaluation is more complex than traditional ML!
	\end{frame}
	
	\begin{frame}{Slide 72: Content Quality Metrics}
		\fontsize{6.5}{7.5}\selectfont
		\textbf{With Reference Materials:}
		\begin{itemize}
			\item Compare output to reference
			\item Check accuracy
			\item May miss qualitative aspects
			\item Limited to available references
		\end{itemize}
		
		\vspace{0.1cm}
		\textbf{Without Reference Materials:}
		\begin{enumerate}
			\item \textbf{Quality-Based}
			\begin{itemize}
				\item Compare reconstructions
				\item LLM-generated reference
			\end{itemize}
			\item \textbf{Entailment-Based}
			\begin{itemize}
				\item Check if text supports hypothesis
				\item Contradiction detection
			\end{itemize}
			\item \textbf{Factuality-Based}
			\begin{itemize}
				\item Verify factual correctness
				\item Check against knowledge bases
			\end{itemize}
		\end{enumerate}
		
		\textcolor{myblue}{\textbf{Key Insight:}} Multiple approaches to assess quality!
	\end{frame}
	
	\begin{frame}{Slide 73: Hallucination}
		\fontsize{6.5}{7.5}\selectfont
		\textbf{What is Hallucination?}
		\begin{itemize}
			\item Output appears reasonable but is factually incorrect
			\item Not supported by input
			\item Nonsensical or improbable scenarios
			\item Caused by probabilistic generation
		\end{itemize}
		
		\vspace{0.1cm}
		\textbf{Why It Happens:}
		\begin{enumerate}
			\item \textbf{Probabilistic Nature}
			\begin{itemize}
				\item Tokens chosen based on probabilities
				\item Higher temperature = more variation
				\item May pick incorrect tokens
			\end{itemize}
			\item \textbf{Knowledge Gaps}
			\begin{itemize}
				\item Model may not have information
				\item Tries to fill gaps
				\item Creates plausible answers
			\end{itemize}
		\end{enumerate}
		
		\vspace{0.1cm}
		\textbf{Importance:}
		\begin{itemize}
			\item Must examine model for hallucination propensity
			\item Critical for risk management applications
			\item Need verification mechanisms
		\end{itemize}
		
		\textcolor{myblue}{\textbf{Key Insight:}} Hallucination is a serious GenAI risk!
	\end{frame}
	
	\begin{frame}{Slide 74: Robustness Evaluation}
		\fontsize{6.5}{7.5}\selectfont
		\textbf{Out of Distribution (OOD) Robustness:}
		\begin{itemize}
			\item Performance on unseen data
			\item Novel inputs
			\item Tests generalization
		\end{itemize}
		
		\vspace{0.1cm}
		\textbf{Adversarial Robustness:}
		\begin{itemize}
			\item Resistance to attacks
			\item Malicious inputs
			\item Designed to cause errors
		\end{itemize}
		
		\vspace{0.1cm}
		\textbf{Jailbreak Resistance:}
		\begin{itemize}
			\item Bypass safety features
			\item Override policy constraints
			\item Test model boundaries
		\end{itemize}
		
		\vspace{0.1cm}
		\textbf{Prompt Sensitivity:}
		\begin{itemize}
			\item Sensitivity to prompt changes
			\item Same question, different phrasing
			\item Stability assessment
		\end{itemize}
		
		\textcolor{myblue}{\textbf{Key Insight:}} Robustness tests model security and reliability!
	\end{frame}
	

\begin{frame}{Slide 75: Bias and Fairness Metrics}
	\fontsize{6.5}{7.5}\selectfont
	
	\textbf{Demographic Parity:}
	\begin{itemize}
		\item Positive predictions are the same across groups
		\item Promotes fair representation
	\end{itemize}
	
	\vspace{0.1cm}
	\textbf{WEAT (Word Embedding Association Test):}
	\begin{itemize}
		\item Measures bias using cosine similarity
		\item Compares associations between sets of words
		\item Produces a normalized score
	\end{itemize}
	
	\vspace{0.1cm}
	\textbf{Gender Bias Score:}
	\begin{itemize}
		\item Indicates gender balance
		\item $> 0$ = male bias
		\item $< 0$ = female bias
	\end{itemize}
	
	\vspace{0.1cm}
	\textbf{Stereotype Score:}
	\begin{itemize}
		\item Measures preference for stereotypical associations
		\item $50$ = neutral
		\item $100$ = strong stereotypical bias
	\end{itemize}
	
	\textcolor{myblue}{\textbf{Key Insight:}}
	Bias metrics help identify fairness issues!
\end{frame}


	
	\begin{frame}{Slide 76: Operational Performance Metrics}
		\fontsize{6.5}{7.5}\selectfont
		\textbf{Time to First Token:}
		\begin{itemize}
			\item Time between request and first token
			\item Perceived responsiveness
			\item Important for user experience
		\end{itemize}
		
		\vspace{0.1cm}
		\textbf{Time Between Tokens:}
		\begin{itemize}
			\item Average interval between tokens
			\item Generation speed
			\item Affects flow
		\end{itemize}
		
		\vspace{0.1cm}
		\textbf{Latency:}
		\begin{itemize}
			\item Response time for query
			\item Overall speed
			\item Critical for real-time applications
		\end{itemize}
		
		\vspace{0.1cm}
		\textbf{Throughput:}
		\begin{itemize}
			\item Number of queries processed
			\item Per unit time
			\item Capacity measure
		\end{itemize}
		
		\textcolor{myblue}{\textbf{Key Insight:}} Speed matters for practical applications!
	\end{frame}
	
	\begin{frame}{Slide 77: Hallucination Metrics}
		\fontsize{6.5}{7.5}\selectfont
		\textbf{Hallucination Rate:}
		\begin{itemize}
			\item Measures model's propensity for inaccurate output
			\item Quantifies hallucination frequency
			\item Essential risk metric
		\end{itemize}
		
		\vspace{0.1cm}
		\textbf{Semantic Entropy:}
		\begin{itemize}
			\item Measures uncertainty/ambiguity
			\item Higher entropy = more uncertainty
			\item Indicates output confidence
		\end{itemize}
		
		\vspace{0.1cm}
		\textbf{How to Measure:}
		\begin{itemize}
			\item Compare to known facts
			\item Human evaluation
			\item Automated fact-checking
			\item Reference-based assessment
		\end{itemize}
		
		\vspace{0.1cm}
		\textbf{Application Example:}
		\begin{itemize}
			\item TruthfulQA benchmark
			\item Tests factual accuracy
			\item Identifies hallucination patterns
		\end{itemize}
		
		\textcolor{myblue}{\textbf{Key Insight:}} Hallucination metrics identify unreliable outputs!
	\end{frame}
	
	\begin{frame}{Slide 78: Text Quality Metrics}
		\fontsize{6.5}{7.5}\selectfont
		\textbf{Exact Match:}
		\begin{itemize}
			\item Binary match/no match
			\item Strictest metric
			\item 0 = no match, 1 = perfect match
		\end{itemize}
		
		\vspace{0.1cm}
		\textbf{BLEU (Bilingual Evaluation Understudy):}
		\begin{itemize}
			\item n-gram precision
			\item Adds brevity penalty
			\item 0-1 scale
		\end{itemize}
		
		\vspace{0.1cm}
		\textbf{ROUGE (Recall-Oriented Understudy):}
		\begin{itemize}
			\item For summaries/translations
			\item n-gram matching
			\item 0-1 scale
		\end{itemize}
		
		\vspace{0.1cm}
		\textbf{BERTScore:}
		\begin{itemize}
			\item BERT sentence embeddings
			\item Compares to reference
			\item Semantic similarity
		\end{itemize}
		
		\textcolor{myblue}{\textbf{Key Insight:}} Multiple metrics for different evaluation needs!
	\end{frame}
	
	\begin{frame}{Slide 79: Image Quality Metrics}
		\fontsize{6.5}{7.5}\selectfont
		\textbf{Inception Score:}
		\begin{itemize}
			\item Classifies generated images
			\item Uses pre-trained Inception-V3
			\item Measures image distinctiveness
			\item Entropy of label distribution
		\end{itemize}
		
		\vspace{0.1cm}
		\textbf{FID (Fréchet Inception Distance):}
		\begin{itemize}
			\item Compares generated vs real images
			\item Distribution parameters
			\item Mean and covariance
			\item Lower FID = better quality
		\end{itemize}
		
		\vspace{0.1cm}
		\textbf{Audio Quality Metrics:}
		\begin{itemize}
			\item Signal to noise ratio
			\item Frequency response
			\item Spectral convergence
		\end{itemize}
		
		\textcolor{myblue}{\textbf{Key Insight:}} Different metrics for different content types!
	\end{frame}
	
	\begin{frame}{Slide 80: Benchmarking LLMs}
		\fontsize{6.5}{7.5}\selectfont
		\textbf{What is Benchmarking?}
		\begin{itemize}
			\item Structured, repeatable tests
			\item Assess abilities across domains
			\item Enable model comparison
			\item Standardized evaluation
		\end{itemize}
		
		\vspace{0.1cm}
		\textbf{Popular Benchmarks:}
		\begin{enumerate}
			\item \textbf{MMLU}
			\begin{itemize}
				\item 57 subjects
				\item Multiple choice questions
				\item Humanities, STEM, social sciences
			\end{itemize}
			\item \textbf{MATH}
			\begin{itemize}
				\item Mathematical capabilities
				\item Problem solving
			\end{itemize}
			\item \textbf{GSM8K}
			\begin{itemize}
				\item Grade school math
				\item Word problems
			\end{itemize}
			\item \textbf{HumanEval}
			\begin{itemize}
				\item Code generation
				\item Programming tasks
			\end{itemize}
		\end{enumerate}
		
		\textcolor{myblue}{\textbf{Key Insight:}} Benchmarks enable fair model comparison!
	\end{frame}
	
	\begin{frame}{Slide 81: MMLU Example Questions}
		\fontsize{6.5}{7.5}\selectfont
		\textbf{World Religions Example:}
		"The Great Cloud Sutra prophesied the imminent arrival of which person?"
		A) Maitreya (Milo) \checkmark
		B) The Buddha
		C) Zhou Dunyi
		D) Wang Yangming
		
		\vspace{0.1cm}
		\textbf{Sociology Example:}
		"Which statement corresponds with differential association theory?"
		A) "If all your friends jumped off a bridge, I suppose you would too." \checkmark
		B) "You should be proud to be a part of this organization."
		C) "If the door is closed, try the window."
		D) "Once a thief, always a thief."
		
		\vspace{0.1cm}
		\textbf{Physics Example:}
		"A model airplane flies slower into the wind and faster with wind at its back. Launched at right angles to the wind, its groundspeed compared with flying in still air is:"
		A) the same
		B) greater \checkmark
		C) less
		D) either greater or less
		
		\textcolor{myblue}{\textbf{Key Insight:}} MMLU tests broad knowledge across subjects!
	\end{frame}
	
	\begin{frame}{Slide 82: LLM Leaderboards}
		\fontsize{6.5}{7.5}\selectfont
		\textbf{What are Leaderboards?}
		\begin{itemize}
			\item Aggregate benchmark results
			\item Show overall model performance
			\item Easy comparison
			\item Frequently updated
		\end{itemize}
		
		\vspace{0.1cm}
		\textbf{Examples:}
		\begin{enumerate}
			\item \textbf{Hugging Face Open LLM Leaderboard}
			\item \textbf{Chatbot Arena}
			\item \textbf{HELM (Holistic Evaluation)}
		\end{enumerate}
		
		\vspace{0.1cm}
		\textbf{HELM Rankings (2025):}
		\begin{tabular}{|c|c|}
			\hline
			\textbf{Model} & \textbf{Mean Win Rate} \\
			\hline
			GPT-4o (2024-05-13) & 0.938 \\
			GPT-4o (2024-08-06) & 0.928 \\
			DeepSeek v3 & 0.908 \\
			Claude 3.5 Sonnet & 0.885 \\
			Amazon Nova Pro & 0.885 \\
			\hline
		\end{tabular}
		
		\textcolor{myblue}{\textbf{Key Insight:}} Leaderboards show model rankings!
	\end{frame}
	
	\begin{frame}{Slide 83: Benchmark Issues}
		\fontsize{6.5}{7.5}\selectfont
		\textbf{Data Leakage (Contamination):}
		\begin{itemize}
			\item Benchmark data in training
			\item Model "memorizes" answers
			\item Inflated performance
			\item False indication of ability
		\end{itemize}
		
		\vspace{0.1cm}
		\textbf{Benchmark Saturation:}
		\begin{itemize}
			\item Models reach maximum scores
			\item Cannot differentiate performance
			\item Benchmarks become outdated
			\item Need new challenges
		\end{itemize}
		
		\vspace{0.1cm}
		\textbf{Lack of Real-World Context:}
		\begin{itemize}
			\item Hypothetical scenarios
			\item May not reflect actual use
			\item Different from real problems
		\end{itemize}
		
		\textcolor{myblue}{\textbf{Key Insight:}} Benchmarks have important limitations!
	\end{frame}
	
	\begin{frame}{Slide 84: Benchmark Issues (continued)}
		\fontsize{6.5}{7.5}\selectfont
		\textbf{Cherry Picking:}
		\begin{itemize}
			\item Only highlight good results
			\item Omit weaker benchmarks
			\item Misleading representation
			\item Selective reporting
		\end{itemize}
		
		\vspace{0.1cm}
		\textbf{Incomplete Comparison:}
		\begin{itemize}
			\item Some leaderboards only show open-source models
			\item Exclude proprietary models
			\item Incomplete landscape
			\item Limited comparison
		\end{itemize}
		
		\vspace{0.1cm}
		\textbf{Impact:}
		\begin{itemize}
			\item Hard to assess true performance
			\item Misleading comparisons
			\item Need critical evaluation
		\end{itemize}
		
		\textcolor{myblue}{\textbf{Key Insight:}} Be critical when interpreting benchmarks!
	\end{frame}
	
	\begin{frame}{Slide 85: GenAI Performance Summary}
		\fontsize{6.5}{7.5}\selectfont
		\textbf{Key Metrics Categories:}
		\begin{enumerate}
			\item \textbf{Content Quality}
			\begin{itemize}
				\item Accuracy, fluency, coherence
				\item BLEU, ROUGE, BERTScore
			\end{itemize}
			\item \textbf{Hallucination}
			\begin{itemize}
				\item Factual accuracy
				\item Semantic entropy
			\end{itemize}
			\item \textbf{Robustness}
			\begin{itemize}
				\item OOD performance
				\item Adversarial resistance
				\item Jailbreak resistance
			\end{itemize}
			\item \textbf{Bias and Fairness}
			\begin{itemize}
				\item Demographic parity
				\item Gender bias
				\item Stereotype scores
			\end{itemize}
			\item \textbf{Operational Performance}
			\begin{itemize}
				\item Latency, throughput, token efficiency
			\end{itemize}
		\end{enumerate}
		
		\textcolor{myblue}{\textbf{Key Insight:}} Comprehensive evaluation requires multiple metrics!
	\end{frame}
	
	% ============================================================
	% SECTION 7: APPLICATIONS - SLIDES 86-100
	% ============================================================
	
	\begin{frame}{Slide 86: GenAI Text Processing Applications}
		\fontsize{6.5}{7.5}\selectfont
		\textbf{Writing Assistance:}
		\begin{itemize}
			\item Report generation
			\item Document summarization
			\item Legal analysis support
			\item Litigation support
		\end{itemize}
		
		\vspace{0.1cm}
		\textbf{Risk Monitoring:}
		\begin{itemize}
			\item Monitor news feeds
			\item Track customer opinions
			\item Social media analysis
			\item Identify potential risks
		\end{itemize}
		
		\vspace{0.1cm}
		\textbf{Market Analysis:}
		\begin{itemize}
			\item Analyze incoming market information
			\item Economic data analysis
			\item News flow analysis
			\item Predict short-term price movements
		\end{itemize}
		
		\textcolor{myblue}{\textbf{Key Insight:}} GenAI automates text-intensive tasks!
	\end{frame}
	
	\begin{frame}{Slide 87: GenAI in Credit and Compliance}
		\fontsize{6.5}{7.5}\selectfont
		\textbf{Credit Underwriting:}
		\begin{itemize}
			\item Process structured and unstructured data
			\item Generate reports
			\item In-depth borrower examination
			\item Reduce decision time
		\end{itemize}
		
		\vspace{0.1cm}
		\textbf{Compliance Monitoring:}
		\begin{itemize}
			\item Analyze SEC filings
			\item Review company reports
			\item Detect emerging risks
			\item Regulatory compliance
		\end{itemize}
		
		\vspace{0.1cm}
		\textbf{Specialized LLMs:}
		\begin{itemize}
			\item Bloomberg GPT
			\item FinGPT
			\item Financial news analysis
			\item Trading signal generation
		\end{itemize}
		
		\textcolor{myblue}{\textbf{Key Insight:}} GenAI transforms credit and compliance work!
	\end{frame}
	
	\begin{frame}{Slide 88: Code Generation}
		\fontsize{6.5}{7.5}\selectfont
		\textbf{Applications:}
		\begin{itemize}
			\item Generate code from descriptions
			\item Debug existing code
			\item Reduce development costs
			\item Accelerate development
		\end{itemize}
		
		\vspace{0.1cm}
		\textbf{Specialist Models:}
		\begin{enumerate}
			\item \textbf{Codex (OpenAI)}
			\begin{itemize}
				\item Code generation
				\item Natural language to code
			\end{itemize}
			\item \textbf{GitHub Copilot (Microsoft)}
			\begin{itemize}
				\item AI pair programmer
				\item Code suggestions
			\end{itemize}
			\item \textbf{Cursor}
			\begin{itemize}
				\item AI-powered IDE
				\item Code assistance
			\end{itemize}
			\item \textbf{Claude Code (Anthropic)}
			\begin{itemize}
				\item Code generation
				\item Analysis and debugging
			\end{itemize}
		\end{enumerate}
		
		\textcolor{myblue}{\textbf{Key Insight:}} GenAI is transforming software development!
	\end{frame}
	
	\begin{frame}{Slide 89: Chatbots and Virtual Assistants}
		\fontsize{6.5}{7.5}\selectfont
		\textbf{Applications:}
		\begin{enumerate}
			\item \textbf{Customer Support}
			\begin{itemize}
				\item Online customer service
				\item 24/7 availability
				\item Quick responses
			\end{itemize}
			\item \textbf{Risk Management}
			\begin{itemize}
				\item Regulatory information
				\item Compliance guidance
				\item Professional support
			\end{itemize}
			\item \textbf{Healthcare}
			\begin{itemize}
				\item Symptom analysis
				\item Medical information
				\item Patient support
			\end{itemize}
			\item \textbf{Investment}
			\begin{itemize}
				\item Portfolio identification
				\item Investment advice
				\item Client support
			\end{itemize}
		\end{enumerate}
		
		\textcolor{myblue}{\textbf{Key Insight:}} Chatbots improve customer and professional services!
	\end{frame}
	
	\begin{frame}{Slide 90: Fraud and Anomaly Detection}
		\fontsize{6.5}{7.5}\selectfont
		\textbf{Capabilities:}
		\begin{itemize}
			\item Automate monitoring and processing
			\item Compile transaction data
			\item Scrutinize across departments
			\item Identify anomalies/outliers
		\end{itemize}
		
		\vspace{0.1cm}
		\textbf{Benefits:}
		\begin{enumerate}
			\item \textbf{Speed}
			\begin{itemize}
				\item Real-time detection
				\item Immediate alerts
			\end{itemize}
			\item \textbf{Comprehensiveness}
			\begin{itemize}
				\item Cross-department analysis
				\item Pattern identification
			\end{itemize}
			\item \textbf{Reduced False Positives}
			\begin{itemize}
				\item Better accuracy
				\item Fewer false alarms
			\end{itemize}
		\end{enumerate}
		
		\vspace{0.1cm}
		\textbf{Example:}
		\begin{itemize}
			\item Mastercard's Decision Intelligence Pro
			\item Dedicated GenAI fraud detection tool
			\item Advanced pattern recognition
		\end{itemize}
		
		\textcolor{myblue}{\textbf{Key Insight:}} GenAI enhances fraud detection capabilities!
	\end{frame}
	
	\begin{frame}{Slide 91: Cybersecurity Applications}
		\fontsize{6.5}{7.5}\selectfont
		\textbf{Threat Monitoring:}
		\begin{itemize}
			\item Monitor emails and data streams
			\item Identify potential cyber threats
			\item Cross-reference with system profile
			\item Pinpoint specific vulnerabilities
		\end{itemize}
		
		\vspace{0.1cm}
		\textbf{Automated Response:}
		\begin{enumerate}
			\item \textbf{Threat Identification}
			\begin{itemize}
				\item Real-time detection
				\item Pattern recognition
				\item Anomaly detection
			\end{itemize}
			\item \textbf{Aleriting}
			\begin{itemize}
				\item Risk managers notified
				\item Real-time alerts
				\item Immediate awareness
			\end{itemize}
			\item \textbf{Solution Sourcing}
			\begin{itemize}
				\item Proactive patch identification
				\item Vendor-approved sources
				\item Coordinated response
			\end{itemize}
		\end{enumerate}
		
		\textcolor{myblue}{\textbf{Key Insight:}} GenAI enables proactive cybersecurity!
	\end{frame}
	
	\begin{frame}{Slide 92: Cybersecurity Examples}
		\fontsize{6.5}{7.5}\selectfont
		\textbf{Cisco Systems:}
		\begin{itemize}
			\item Custom LLM for malware detection
			\item Security analysis
			\item Threat identification
		\end{itemize}
		
		\vspace{0.1cm}
		\textbf{CrowdStrike:}
		\begin{itemize}
			\item Charlotte AI for cybersecurity
			\item Automated threat response
			\item Security operations
		\end{itemize}
		
		\vspace{0.1cm}
		\textbf{Key Features:}
		\begin{enumerate}
			\item \textbf{Automated Analysis}
			\begin{itemize}
				\item Processes large data volumes
				\item Identifies patterns
				\item Reduces response time
			\end{itemize}
			\item \textbf{Proactive Defense}
			\begin{itemize}
				\item Pre-emptive identification
				\item Automated patching
				\item Coordinated response
			\end{itemize}
		\end{enumerate}
		
		\textcolor{myblue}{\textbf{Key Insight:}} LLMs strengthen cybersecurity defenses!
	\end{frame}
	
	\begin{frame}{Slide 93: Image Generation Applications}
		\fontsize{6.5}{7.5}\selectfont
		\textbf{Rapid Prototyping:}
		\begin{itemize}
			\item Generate images from text
			\item Quick concept visualization
			\item Iterative design
		\end{itemize}
		
		\vspace{0.1cm}
		\textbf{Applications:}
		\begin{enumerate}
			\item \textbf{Marketing Materials}
			\begin{itemize}
				\item Campaign visuals
				\item Advertisements
				\item Social media content
			\end{itemize}
			\item \textbf{Book Covers}
			\begin{itemize}
				\item Custom artwork
				\item Quick generation
			\end{itemize}
			\item \textbf{Illustrations}
			\begin{itemize}
				\item Presentation graphics
				\item Reports and documents
				\item Educational materials
			\end{itemize}
		\end{enumerate}
		
		\vspace{0.1cm}
		\textbf{Key Models:}
		\begin{itemize}
			\item DALL-E (OpenAI)
			\item Stable Diffusion
			\item Midjourney
		\end{itemize}
		
		\textcolor{myblue}{\textbf{Key Insight:}} Image generation accelerates creative workflows!
	\end{frame}
	
	\begin{frame}{Slide 94: Text, Speech, Video Conversion}
		\fontsize{6.5}{7.5}\selectfont
		\textbf{Multimodal Capabilities:}
		\begin{enumerate}
			\item \textbf{Text to Speech}
			\begin{itemize}
				\item OpenAI Whisper
				\item Text-to-Speech (TTS)
				\item Voice synthesis
			\end{itemize}
			\item \textbf{Speech to Text}
			\begin{itemize}
				\item Transcription
				\item Meeting notes
				\item Accessibility
			\end{itemize}
			\item \textbf{Text to Video}
			\begin{itemize}
				\item OpenAI's Sora
				\item Google Gemini
				\item Realistic video creation
			\end{itemize}
		\end{enumerate}
		
		\vspace{0.1cm}
		\textbf{Applications:}
		\begin{itemize}
			\item Content creation
			\item Accessibility services
			\item Educational materials
			\item Marketing and advertising
		\end{itemize}
		
		\textcolor{myblue}{\textbf{Key Insight:}} Multimodal models convert between content types!
	\end{frame}
	
	\begin{frame}{Slide 95: Autonomous Driving}
		\fontsize{6.5}{7.5}\selectfont
		\textbf{Development:}
		\begin{itemize}
			\item Under development since 1970s
			\item AI is crucial component
			\item Semi-autonomous systems deployed
		\end{itemize}
		
		\vspace{0.1cm}
		\textbf{Current Systems:}
		\begin{enumerate}
			\item \textbf{Waymo (USA)}
			\begin{itemize}
				\item Advanced autonomous driving
				\item Public deployment
			\end{itemize}
			\item \textbf{Pony.ai (China)}
			\begin{itemize}
				\item Semi-autonomous systems
				\item Urban deployment
			\end{itemize}
		\end{enumerate}
		
		\vspace{0.1cm}
		\textbf{Future Outlook:}
		\begin{itemize}
			\item More robust systems expected
			\item Increased adoption
			\item Continuous improvement
			\item Safety remains concern
		\end{itemize}
		
		\textcolor{myblue}{\textbf{Key Insight:}} AI is essential for autonomous driving!
	\end{frame}
	
	\begin{frame}{Slide 96: Insurance LLM Use Case}
		\fontsize{6.5}{7.5}\selectfont
		\textbf{EXL Insurance LLM:}
		\begin{itemize}
			\item Built with NVIDIA AI Enterprise
			\item Insurance-specific LLM
			\item Customized with insurance data
			\item Outperforms generic LLMs
		\end{itemize}
		
		\vspace{0.1cm}
		\textbf{Applications:}
		\begin{enumerate}
			\item \textbf{Claims Processing}
			\begin{itemize}
				\item Streamline document review
				\item Faster processing
				\item Improved accuracy
			\end{itemize}
			\item \textbf{Underwriting}
			\begin{itemize}
				\item Better data interpretation
				\item Informed decisions
				\item Reduced workload
			\end{itemize}
			\item \textbf{Customer Service}
			\begin{itemize}
				\item Automated support
				\item Quick responses
				\item Consistent service
			\end{itemize}
		\end{enumerate}
		
		\textcolor{myblue}{\textbf{Key Insight:}} Domain-specific LLMs outperform general models!
	\end{frame}
	
	\begin{frame}{Slide 97: EXL Insurance LLM Results}
		\fontsize{6.5}{7.5}\selectfont
		\textbf{Performance Improvements:}
		\begin{enumerate}
			\item \textbf{30\% Reduction in Operational Costs}
			\begin{itemize}
				\item More efficient claims processing
				\item Streamlined underwriting
				\item Automated customer service
			\end{itemize}
			\item \textbf{Improved Accuracy}
			\begin{itemize}
				\item Better data interpretation
				\item Enhanced decision-making
				\item Reduced errors
			\end{itemize}
			\item \textbf{Reduced Manual Workload}
			\begin{itemize}
				\item Staff focus on complex tasks
				\item Increased productivity
				\item Better resource allocation
			\end{itemize}
		\end{enumerate}
		
		\vspace{0.1cm}
		\textbf{Key Features:}
		\begin{itemize}
			\item Secure infrastructure
			\item Regular audits
			\item Regulatory compliance
			\item Enterprise-ready
		\end{itemize}
		
		\textcolor{myblue}{\textbf{Key Insight:}} Insurance LLMs deliver significant business value!
	\end{frame}
	
	\begin{frame}{Slide 98: Investment Firms Using LLMs}
		\fontsize{6.5}{7.5}\selectfont
		\textbf{BlackRock:}
		\begin{itemize}
			\item Using LLMs in investment process
			\item Market analysis
			\item Investment decisions
			\item Portfolio management
		\end{itemize}
		
		\vspace{0.1cm}
		\textbf{Bridgewater:}
		\begin{itemize}
			\item \$2 billion fund run by machine learning
			\item Advanced analytics
			\item Quantitative strategies
			\item AI-driven decisions
		\end{itemize}
		
		\vspace{0.1cm}
		\textbf{AQR:}
		\begin{itemize}
			\item LLM integration
			\item Investment strategies
			\item Risk management
			\item Data analysis
		\end{itemize}
		
		\textcolor{myblue}{\textbf{Key Insight:}} Leading firms are adopting LLMs for investment!
	\end{frame}
	
	\begin{frame}{Slide 99: GenAI Impact Summary}
		\fontsize{6.5}{7.5}\selectfont
		\textbf{Key Areas of Impact:}
		\begin{enumerate}
			\item \textbf{Fraud Detection}
			\begin{itemize}
				\item Automated monitoring
				\item Pattern recognition
				\item Real-time alerts
			\end{itemize}
			\item \textbf{Customer Service}
			\begin{itemize}
				\item Chatbots and virtual assistants
				\item 24/7 availability
				\item Improved response quality
			\end{itemize}
			\item \textbf{Document Processing}
			\begin{itemize}
				\item Summarization
				\item Analysis
				\item Compliance
			\end{itemize}
			\item \textbf{Investment}
			\begin{itemize}
				\item Market analysis
				\item Trading signals
				\item Portfolio management
			\end{itemize}
		\end{enumerate}
		
		\textcolor{myblue}{\textbf{Key Insight:}} GenAI is transforming every financial sector!
	\end{frame}
	

\begin{frame}{Slide 100: Final Summary}
	\fontsize{6.5}{7.5}\selectfont
	
	\textbf{\textcolor{myblue}{Generative AI Overview:}}
	\begin{itemize}
		\item Creates new content (text, images, audio, and video)
		\item Rapidly evolving technology
		\item Broad range of applications
	\end{itemize}
	
	\vspace{0.1cm}
	\textbf{\textcolor{myblue}{Key Technologies:}}
	\begin{itemize}
		\item Word2Vec: Creates word embeddings
		\item RNNs/LSTMs: Sequential data processing
		\item Transformers: Attention-based, parallel processing
		\item LLMs: Large-scale, pre-trained models
	\end{itemize}
	
	\vspace{0.1cm}
	\textbf{\textcolor{myblue}{Performance Evaluation:}}
	\begin{itemize}
		\item Content quality metrics
		\item Hallucination detection
		\item Robustness and fairness
		\item Operational performance
	\end{itemize}
	
	\vspace{0.1cm}
	\textbf{\textcolor{myblue}{Applications in Finance:}}
	\begin{itemize}
		\item Fraud detection, customer service, and document analysis
		\item Risk management, compliance, and investment
		\item Credit underwriting and cybersecurity
	\end{itemize}
	
	\textcolor{myred}{\textbf{Exam Success:}}
	Understand the technology, evaluation, and applications!
\end{frame}


	
\end{document}